\section{AIME I 2018}
        \begin{enumerate}
        	\item (\textit{AIME I 2018}) Let $S$ be the number of ordered pairs of integers $(a,b)$ with $1 \leq a \leq 100$ and $b \geq 0$ such that the polynomial $x^2+ax+b$ can be factored into the product of two (not necessarily distinct) linear factors with integer coefficients. Find the remainder when $S$ is divided by $1000$.


 
 \textbf{Answer: }600



\textbf{Solution 1}

Let the linear factors be $(x+c)(x+d)$.

 Then, $a=c+d$ and $b=cd$.

 We know that $1\le a\le 100$ and $b\ge 0$, so $c$ and $d$ both have to be non-negative

 However, $a$ cannot be $0$, so at least one of $c$ and $d$ must be greater than $0$, ie positive.

 Also, $a$ cannot be greater than $100$, so $c+d$ must be less than or equal to $100$.

 Essentially, if we plot the solutions, we get a triangle on the coordinate plane with vertices $(0,0), (0, 100),$ and $(100,0)$. Remember that $(0,0)$ does not work, so there is a square with the top right corner $(1,1)$.

 Note that $c$ and $d$ are interchangeable since they end up as $a$ and $b$ in the end anyways. Thus, we simply draw a line from $(1,1)$ to $(50,50)$, designating one of the halves as our solution (since the other side is simply the coordinates flipped).

 We note that the pattern from $(1,1)$ to $(50,50)$ is $2+3+4+\dots+51$ solutions and from $(51, 49)$ to $(100,0)$ is $50+49+48+\dots+1$ solutions, since we can decrease the $y$-value by $1$ until $0$ for each coordinate.

 Adding up gives 
 \[\dfrac{2+51}{2}\cdot 50+\dfrac{50+1}{2}\cdot 50.\]This gives us $2600$, and $2600\equiv 600 \bmod{1000}.$

 Thus, the answer is: 
 \[\boxed{600}.\]




\textbf{Solution 2}

Similar to the previous solution, we plot the triangle and cut it in half. Then, we find the number of boundary points, which is $101+51+51-3$, or just $200$. Using Pick's theorem, we know that the area of the half-triangle, which is $2500$, is just $I+100-1$. That means that there are $2401$ interior points, plus $200$ boundary points, which is $2601$. However, $(0,0)$ does not work, so the answer is 
 \[\boxed{600}.\]




\textbf{Solution 3 (less complicated)}

Notice that for $x^2+ax+b$ to be true, for every $a$, $b$ will always be the product of the possibilities of how to add two integers to $a$. For example, if $a=3$, $b$ will be the product of $(3,0)$ and $(2,1)$, as those two sets are the only possibilities of adding two integers to $a$. Note that order does not matter. If we just do some simple casework, we find out that: 

 if $a$ is odd, there will always be $\left\lceil\frac{a}{2}\right\rceil$ $\left(\text{which is also }\frac{a+1}{2}\right)$ possibilities of adding two integers to $a$.

 if $a$ is even, there will always be $\frac{a}{2}+1$ possibilities of adding two integers to $a$. 

 Using the casework, we have $1+2+2+3+3+...50+50+51$ possibilities. This will mean that the answer is 
 \[\frac{(1+51)\cdot100}{2}\Rightarrow52\cdot50=2600\] possibilities.

 Thus, our solution is $2600\bmod {1000}\equiv\boxed{600}$.

 Solution by IronicNinja~

 



\textbf{Solution 4}

Let's write the linear factors as $(x+n)(x+m)$.

 Then we can write them as: $a=n+m, b=nm$.

 $n$ or $m$ has to be a positive integer as a cannot be 0.

 $n+m$ has to be between $1$ and $100$, as a cannot be over $100$.

 Excluding $a=1$, we can see there is always a pair of $2$ a-values for a certain amount of b-values. 

 For instance, $a=2$ and $a=3$ both have $2$ b-values. $a=4$ and $a=5$ both have $3$ b-values.

 We notice the pattern of the number of b-values in relation to the a-values:$1, 2, 2, 3, 3, 4, 4\ldots$

 \textbf{Ending 1} We see the pattern $1, 2; 2, 3; 3, 4; ...$. There are 50 pairs of $i, i+1$ in this pattern, and each pair sums to $2i+1$. So the pattern condenses to $3, 5, 7, ...$ for 50 terms. This is just $1+3+5+...$ for 51 terms, minus $1$, or $51^2-1=2601-1=2600\implies\boxed{600}$.

 \textbf{Ending 2} The following link is the URL to the graph I drew showing the relationship between a-values and b-values\href{https://artofproblemsolving.comhttp://artofproblemsolving.com/wiki/index.php?title=File:Screen\_Shot\_2018-04-30\_at\_8.15.00\_PM.png\#file}{http://artofproblemsolving.com/wiki/index.php?title=File:Screen\_Shot\_2018-04-30\_at\_8.15.00\_PM.png\#file}

 
                    \begin{figure}[H]
                        \centering
                        \includegraphics[width=0.5\textwidth]{Solution/AIME_I_2018_1_Solution_Figure_1.png}
                    \end{figure}
                    The pattern continues until $a=100$, and in total, there are $49$ pairs of a-value with the same amount of b-values. The two lone a-values without a pair are, the ($a=1$, amount of b-values=1) in the beginning, and ($a=100$, amount of b-values=51) in the end. 

 Then, we add numbers from the opposite ends of the spectrum, and quickly notice that there are $50$ pairs each with a sum of $52$. $52\cdot50$ gives $2600$ ordered pairs:

 $1+51, 2+50, 2+50, 3+49, 3+49, 4+48, 4+48\ldots$

 When divided by $1000$, it gives the remainder $\boxed{600}$, the answer.

 Solution provided by- YonglaoSlightly modified by Afly

 
Remark : This solution works because no distinct $(a,b), (c,d)$ exist such that $a+b=c+d,ab=cd$ unless $a = d$ and $b = c$




\textbf{Solution 5}

Let's say that the quadratic $x^2 + ax + b$ can be factored into $(x+c)(x+d)$ where $c$ and $d$ are non-negative numbers. We can't have both of them zero because $a$ would not be within bounds. Also, $c+d \leq 100$. Assume that $c < d$. $d$ can be written as $c + x$ where $x \geq 0$. Therefore, $c + d = 2c + x \leq 100$. To find the amount of ordered pairs, we must consider how many values of $x$ are possible for each value of $c$. The amount of possible values of $x$ is given by $100 - 2c + 1$. The $+1$ is the case where $c = d$. We don't include the case where $c = d = 0$, so we must subtract a case from our total. The amount of ordered pairs of $(a,b)$ is:
 \[\left(\sum_{c=0}^{50} (100 - 2c + 1)\right) - 1\]This is an arithmetic progression. 
 \[\frac{(101 + 1)(51)}{2} - 1 = 2601 - 1 = 2600\] When divided by $1000$, it gives the remainder $\boxed{600}$




\textbf{Solution 6(simple)}

By Vietas, the sum of the roots is $-a$ and the product is $b$. Therefore, both roots are nonpositive. For each value of $a$ from $1$ to $100$, the number of $b$ values is the number of ways to sum two numbers between $0$ and $a-1$ inclusive to $a$. This is just $1 + 2 + 2 + 3 + 3 +... 50 + 50 + 51 = 2600$. Thus, the answer is $\boxed{600}$

 -bron jiang

 



\textbf{Solution 7 (less room for error)}

Similar to solution 1 we plot the triangle and half it. From dividing the triangle in half we are removing the other half of answers that are just flipped coordinates. We notice that we can measure the length of the longest side of the half triangle which is just from $0$ to $100$, so the number of points on that line is is $101$. The next row has length $99$, the one after that has length $97$, and so on. We simply add this arithmetic series of odd integers $101+99+97+...+1$. This is 
 \[\frac{(101+1)(51)}{2} = 2601\] Or you can notice that this is the sum of the first $51$ odd terms, which is just $51^2 = 2601$.However, $(0,0)$ is the singular coordinate that does not work, so the answer is $(2601-1)\bmod {1000}\equiv\boxed{600}$

 Solution by Damian Kim~

 



\textbf{Solution 8 (Sticks and stones)}

Letting $-r$ and $-q$ be the roots, we must find the number of unordered pairs $(r,q)$ such that $1 \le r+q \le 100.$ To do this we define three boxes: $r, q,$ and $t$ for "trash." For example, if we had $t=77,$ we would have any arrangement of $r$ and $q$ summing to $23.$ Note that we cannot have $t=100,$ subtracting one from our total. By stars and bars, we have that the number of ordered pairs $(r,q)$ is ${102 \choose 2} - 1 = 5150.$ However, we are not done: we have double-counted cases such as $(1,2)$ and $(2,1)$ but not cases such as $(4,4).$ Thus our answer will be $\frac{5150 - 50}{2} + 50 \pmod {100} = \boxed{600}.$~ab2024

 



\textbf{Solution 9 (Official MAA)}

The factoring condition is equivalent to the discriminant $a^2-4b$ being equal to $c^2$ for some integer $c.$ Because $b\ge 0,$ the equation $4b=(a-c)(a+c)$ shows that the existence of such a $b$ is equivalent to $a\equiv c\pmod 2$ with $0\le c\le a.$ Thus the number of ordered pairs is 
 \[S=\sum_{a=1}^{100}\left\lceil\frac{a+1}{2}\right\rceil=2600.\] The requested remainder is $600.$




\textbf{Solution 10 (Extra solution that is very very similar to the other solutions)}

First I thought the special case of the problem is when the factors are repeated, which looks like this: 

 \begin{verbatim}
$(x+r)^2$
\end{verbatim}

There are 50 ways for this to happen because r can be from $1$ to $50$ ($(x+50)^2 = x^2+100x+50^2$ and $a\leq 50$).

 Now I thought how can you make the other cases? It looks like $(x-r)(x-s)$ so by Vieta's $r+s=a$. This looks like Stars and Bars

 because $a=1+1+1+1+...$ and you're giving the indistinguishable ones to $r$ and $s$. For $a=1$, there are $\binom{1+2-1}{2-1} = 2$

 ways. For $a=2$, there are $\binom{2+2-1}{2-1} = 3$ ways. Now you see the pattern which is $2+3+...+100+101$ for $a=1, 2, ...100$. 

 This equals $103\cdot50$ ways in total. However, now you realize that $r$ and $s$ can be flipped and it gives the same. So, 

 subtract the cases when $r = s$ which we found to be $50$. Now from the remaining double-counted $102\cdot50$ divide by two to

 eliminate the double-counting. We have $51\cdot50$, but remember to add back in the $50$ ways that weren't double-counted to begin

 with. That gives us $52\cdot50 = 2600$ ways and the requested remainder is $600.$

 -unhappyfarmer

 


	\item (\textit{AIME I 2018}) The number $n$ can be written in base $14$ as $\underline{a}\text{ }\underline{b}\text{ }\underline{c}$, can be written in base $15$ as $\underline{a}\text{ }\underline{c}\text{ }\underline{b}$, and can be written in base $6$ as $\underline{a}\text{ }\underline{c}\text{ }\underline{a}\text{ }\underline{c}\text{ }$, where $a > 0$. Find the base-$10$ representation of $n$.


 
 \textbf{Answer: }925



\textbf{Solution 1}

We have these equations:$196a+14b+c=225a+15c+b=222a+37c$.Taking the last two we get $3a+b=22c$. Because $c \neq 0$ otherwise $a \ngtr 0$, and $a \leq 5$, $c=1$.

 Then we know $3a+b=22$.Taking the first two equations we see that $29a+14c=13b$. Combining the two gives $a=4, b=10, c=1$. Then we see that $222 \times 4+37 \times1=\boxed{925}$.

 



\textbf{Solution 2}

We know that $196a+14b+c=225a+15c+b=222a+37c$. Combining the first and third equations give that $196a+14b+c=222a+37c$, or 
 \[7b=13a+18c\]The second and third gives $222a+37c=225a+15c+b$, or 
 \[22c-3a=b\]
 \[154c-21a=7b=13a+18c\]
 \[4c=a\]We can have $a=4,8,12,16,20$, but only $a=4$ falls within the possible digits of base $6$. Thus $a=4$, $c=1$, and thus you can find $b$ which equals $10$. Thus, our answer is $4\cdot225+1\cdot15+10=\boxed{925}$.~SHEN KISLAY KAI 2023

 



\textbf{Solution 3 (Official MAA)}

The problem is equivalent to finding a solution to the system of Diophantine equations $196a+14b+c=225a+15c+b$ and $225a+15c+b=216a+36c+6a+c,$ where $1\le a\le 5,\,0\le b\le 13,$ and $0\le c\le 5.$ Simplifying the second equation gives $b=22c-3a.$ Substituting for $b$ in the first equation and simplifying then gives $a=4c,$ so $a = 4$ and $c = 1,$ and the base-$10$ representation of $n$ is $222 \cdot 4 + 37 \cdot 1 = 925.$ It may be verified that $b=10\le 13.$




\textbf{Solution 4 (Simple Modular Arithmetic)}

We're given that $196a+14b+c=225a+15c+b=222a+37c.$ By taking the difference of the first $2$ equalities, we receive $29a+14c=13b.$ Taking $\pmod{13}$, we receive $3a+c \equiv 0 \pmod{13}.$ We receive the following cases: $(a,c)=(4,1)$ or $(3,4).$ (Note that $(2,7)$ doesn't work since $a,c<6$ by third condition). We can just check these two, and find that $(a,b,c)=(4,10,1),$ and just plugging in $(a,c)$ into the third expression we receive $888+37=\boxed{925}.$



	\item (\textit{AIME I 2018}) Kathy has $5$ red cards and $5$ green cards. She shuffles the $10$ cards and lays out $5$ of the cards in a row in a random order. She will be happy if and only if all the red cards laid out are adjacent and all the green cards laid out are adjacent. For example, card orders $RRGGG$, $GGGGR$, or $RRRRR$ will make Kathy happy, but $RRRGR$ will not. The probability that Kathy will be happy is $\frac{m}{n}$, where $m$ and $n$ are relatively prime positive integers. Find $m + n$ .


 
 \textbf{Answer: }157



\textbf{Solution 1}

We have $2+4\cdot 2$ cases total.

 The two are all red and all green. Then, you have 4 of one, 1 of other. 3 of one, 2 of other. 2 of one, 3 of other. 1 of one, 4 of other. Then flip the order, so times two.

 Obviously the denominator is $10\cdot 9\cdot 8\cdot 7\cdot 6$, since we are choosing a card without replacement.

 Then, we have for the numerator for the two of all red and green: 
 \[5\cdot 4\cdot 3\cdot 2\cdot 1.\]

 For the 4 and 1, we have: 
 \[5\cdot 4\cdot 3\cdot 2\cdot 5.\]

 For the 3 and 2, we have: 
 \[5\cdot 4\cdot 3\cdot 5\cdot 4.\]

 For the 2 and 3, we have: 
 \[5\cdot 4\cdot 5\cdot 4\cdot 3.\]

 For the 1 and 4, we have: 
 \[5\cdot 5\cdot 4\cdot 3\cdot 2.\]

 Adding up and remembering to double all of them, since they can be reversed and the 5's can be red or green, we get, after simplifying: $\dfrac{31}{126}$

 Thus the answer is $31 + 126 = \boxed{157}$.-gorefeebuddie

 



\textbf{Solution 2}

Our probability will be $\dfrac{\text{number of "happy" configurations of cards}}{\text{total number of configurations of cards}}.$

 First of all, we have $10$ choices for the first card, $9$ choices for the second card, $8$ choices for the third card, $7$ choices for the fourth card, and $6$ choices for the last card. This gives a total of $10\cdot 9\cdot 8\cdot 7\cdot 6$ possible ways for five cards to be chosen.

 Finding the number of configurations that make Kathy happy is a more difficult task, however, and we will resort to casework to do it.

 First, let's look at the appearances of the "happy configurations" that Kathy likes. Based on the premise of the problem, we can realize that there are ten cases for the appearance of the configurations: 
 \[\text{RRRRR},\] 
 \[\text{GGGGG},\] 
 \[\text{RRRRG},\] 
 \[\text{GGGGR},\] 
 \[\text{RRRGG},\] 
 \[\text{GGGRR},\] 
 \[\text{RRGGG},\] 
 \[\text{GGRRR},\] 
 \[\text{RGGGG},\] 
 \[\text{GRRRR}.\] But this doesn't mean there are 10 "happy configurations" in total-- remember that we've been treating these cards as distinguishable, so we must continue to do so.

 To lighten the load of 10 cases on the human brain, we can note that in the eyes of what we will soon do, $RRRRR$ and $GGGGG$ are effectively equivalent, and therefore may be treated in the same case. We will have to multiply by 2 at the end, though.

 Similarly, we can equate $RRRRG,$ $GGGGR,$ $RGGGG,$ and $GRRRR,$ as well as $RRRGG,$ $GGGRR,$ $RRGGG,$ and $GGRRR,$ so that we just have three cases. We can approach each of these cases with constructive counting.

 Case 1: $RRRRR$-type.

 For this case, there are $5$ available choices for the first card, $4$ available choices for the second card, $3$ for the third card, $2$ for the fourth card, and $1$ for the last card. This leads to a total of $5\cdot 4\cdot 3\cdot 2\cdot 1=120$ configurations for this case. There are $2$ cases of this type.

 Case 2: $RRRRG$-type.

 For this case, there are $5$ available choices for the first card, $4$ available choices for the second card, $3$ for the third card, $2$ for the fourth card, and $5$ choices for the last card (not $1$, because we're doing a new color). This leads to a total of $5\cdot 4\cdot 3\cdot 2\cdot 5=600$ configurations for this case. There are $4$ cases of this type.

 Case 3: $RRRGG$-type.

 For this case, there are $5$ available choices for the first card, $4$ available choices for the second card, $3$ for the third card, $5$ for the fourth card, and $4$ choices for the last card. This leads to a total of $5\cdot 4\cdot 3\cdot 5\cdot 4=1200$ configurations for this case. There are $4$ cases of this type.

 Adding the cases up gives $2\cdot 120+4\cdot 600+4\cdot 1200=7440$ "happy" configurations in total.

 This means that the probability that Kathy is happy will be $\dfrac{7440}{10\cdot 9\cdot 8\cdot 7\cdot 6},$ which simplifies to $\dfrac{31}{126}.$

 So the answer is $31+126=\boxed{157}$




\textbf{Solution 3}

As the problem states, some examples of valid are $RRGGG$, $GGGGR$, and $RRRRR$. Let's use each of these as more general cases. 

 Let $RRGGG$ be the case when there are 2 adjacents of one color, and 3 adjacents of the other color. This yields $4$ combinations ($RRGGG$, $GGRRR$, $RRRGG$, and $GGGRR$). The probability of each of these is equal, equating to $\frac{5}{10}\cdot \frac{4}{9}\cdot \frac{5}{8}\cdot \frac{4}{7}\cdot \frac{3}{6}=\frac{5}{126}$, and since there are $4$ combinations, the probability of this case is $4\cdot \frac{5}{126}=\frac{10}{63}$

 Next case is $GGGGR$. Let this be when there are 4 adjacents of one color, and 1 individual color. Once again, this yields $4$ combinations ($GGGGR$, $RRRRG$, $RGGGG$, and $GRRRR$). The probability of each is the same, equating to $\frac{5}{10}\cdot \frac{4}{9}\cdot \frac{3}{8}\cdot \frac{2}{7}\cdot \frac{5}{6}=\frac{5}{252}$, and since there are $4$ combinations, the probability of this case is $4\cdot \frac{5}{252}=\frac{5}{63}$

 The final case is $RRRRR$, in which there is just an adjacent block of $5$ colors. There are only $2$ combinations this time, each equating to the probability $\frac{5}{10}\cdot \frac{4}{9}\cdot \frac{3}{8}\cdot \frac{2}{7}\cdot \frac{1}{6}=\frac{1}{252}$, and since there are $2$ combinations, the probability of this case is $2\cdot \frac{1}{252}=\frac{1}{126}$. 

 Thus, the total probability is $\frac{10}{63}+\frac{5}{63}+\frac{1}{126}=\frac{31}{126} \implies m+n=\boxed{157}$




\textbf{Solution 4}

Kathy will draw the 10 cards in some order, but the restriction of having all greens in a row and all reds in a row only applies to the first 5 cards drawn. The total number of ways the 10 cards can be drawn is simply 10 choose 5 which is 252. Now we just count the number of possible successful configurations of the ten cards. The first 5 cards can start either be $GRRRR$, $GGRRR$, $GGGRR$, $GGGGR$, $GGGGG$ or the same thing except starting with a red. The number of ways to order $GRRRR$ is the number of ways to order the last 5 cards, which is 5C1. Doing all of the other cases, the total is $(\binom{5}{1}+\binom{5}{2}+\binom{5}{3}+\binom{5}{4}+\binom{5}{5})*2 = 62$. $\frac{62}{252} = \frac{31}{126},$ so the solution is $31 + 126 = \boxed{157}$

 -bradleyguo
Minor LaTeX edits by fasterthanlight

 



\textbf{Solution 5}

Suppose $k$ of the first $5$ cards are red. There are $\binom{5}{k}$ ways to order the last five cards, and either $1$ or $2$ ways to order the first five cards: $1$ if $k=0$ or $k=5$, otherwise $2$. So, the number of orderings that work is
 \[\binom{5}{0} + 2\binom{5}{1} + 2\binom{5}{2} + 2\binom{5}{3} + 2\binom{5}{4} + \binom{5}{5} = 2\sum_{k=0}^{5}\binom{5}{k} - 2 = 2^6 - 2.\]

 The total number of orderings is $\binom{10}{5} = 252,$ so the answer is $\frac{62}{252} = \boxed{\frac{31}{126}}$




\textbf{Solution 6 (Official MAA)}

Assume without loss of generality that the first card laid out is red. Then the arrangements that satisfy Kathy’s requirements are RRRRR, RRRRG, RRRGG, RRGGG, and RGGGG. The probability that Kathy will lay out one of thesearrangements is 
 \[\frac49\cdot\frac38\cdot\frac27\cdot\frac16\] 
 \[\frac49\cdot\frac38\cdot\frac27\cdot\frac56\] 
 \[\frac49\cdot\frac38\cdot\frac57\cdot\frac46\] 
 \[\frac49\cdot\frac58\cdot\frac47\cdot\frac36\] 
 \[+\frac59\cdot\frac48\cdot\frac37\cdot\frac26\] 
 \[\overline{..........................}\] 
 \[\frac{31}{126}\] The requested sum is $31+126=\boxed{157}.$



	\item (\textit{AIME I 2018}) In $\triangle ABC, AB = AC = 10$ and $BC = 12$. Point $D$ lies strictly between $A$ and $B$ on $\overline{AB}$ and point $E$ lies strictly between $A$ and $C$ on $\overline{AC}$ so that $AD = DE = EC$. Then $AD$ can be expressed in the form $\dfrac{p}{q}$, where $p$ and $q$ are relatively prime positive integers. Find $p+q$.


 
 \textbf{Answer: }289



\textbf{Solution (Easiest Law of Cosines)}

We apply Law of Cosines on $\angle A$ twice (one from $\triangle ADE$ and one from $\triangle ABC$),

 \begin{align*}12^2 &= 10^2 + 10^2 - 2(10)(10) \cdot \cos{A} \\[5pt]x^2&=x^2+(10-x)^2-2(x)(10-x)\cdot \cos{A} \end{align*}


 Solving for $\cos{A}$ in both equations, we get\begin{align*}\cos{A} &= \frac{7}{25} \\\cos{A} &= \frac{(10-x)^2}{(2)(10-x)(x)} = \frac{10-x}{2x}\end{align*}
Setting the two equal,\begin{align*} \frac{10-x}{2x} &= \frac{7}{25} \\[5pt]250-25x &= 14x \\[5pt]x &= \frac{250}{39}.\end{align*}
Therefore, our answer is $\boxed{289}$

 Note that this strategy of applying LOC on congruent or supplementary angles is very common. It's also how Stewart's Theorem is derived. 

 \textbf{-RootThreeOverTwo, edits by epiconan}

 



\textbf{Solution 1 (No Trig)}


                    \begin{figure}[H]
                        \centering
                        \includegraphics[width=0.5\textwidth]{Solution/AIME_I_2018_4_Solution_Figure_1.png}
                    \end{figure}
                    

 

 We draw the altitude from $B$ to $\overline{AC}$ to get point $F$. We notice that the triangle's height from $A$ to $\overline{BC}$ is 8 because it is a $3-4-5$ Right Triangle. To find the length of $\overline{BF}$, we let $h$ represent $\overline{BF}$ and set up an equation by finding two ways to express the area. The equation is $(8)(12)=(10)(h)$, which leaves us with $h=9.6$. We then solve for the length $\overline{AF}$, which is done through pythagorean theorm and get $\overline{AF}$ = $2.8$. We can now see that $\triangle AFB$ is a $7-24-25$ Right Triangle. Thus, we set $\overline{AG}$ as $5-$$\frac{x}{2}$, and yield that $\overline{AD}$ $=$ $($ $5-$ $\frac{x}{2}$ $)$ $($ $\frac{25}{7}$ $)$. Now, we can see $x$ = $($ $5-$ $\frac{x}{2}$ $)$ $($ $\frac{25}{7}$ $)$. Solving this equation, we yield $39x=250$, or $x=$ $\frac{250}{39}$. Thus, our final answer is $250+39=\boxed{289}$.

 Diagram edited by Afly

 



\textbf{Solution 2 (Easy Similar Triangles)}

We start by adding a few points to the diagram. Call $F$ the midpoint of $AE$, and $G$ the midpoint of $BC$. (Note that $DF$ and $AG$ are altitudes of their respective triangles). We also call $\angle BAC = \theta$. Since triangle $ADE$ is isosceles, $\angle AED = \theta$, and $\angle ADF = \angle EDF = 90 - \theta$. Since $\angle DEA = \theta$, $\angle DEC = 180 - \theta$ and $\angle EDC = \angle ECD = \frac{\theta}{2}$. Since $FDC$ is a right triangle, $\angle FDC = 180 - 90 - \frac{\theta}{2} = \frac{180-m}{2}$. 

 Since $\angle BAG = \frac{\theta}{2}$ and $\angle ABG = \frac{180-m}{2}$, triangles $ABG$ and $CDF$ are similar by Angle-Angle similarity. Using similar triangle ratios, we have $\frac{AG}{BG} = \frac{CF}{DF}$. $AG = 8$ and $BG = 6$ because there are $2$ $6-8-10$ triangles in the problem. Call $AD = x$. Then $CE = x$, $AE = 10-x$, and $EF = \frac{10-x}{2}$. Thus $CF = x + \frac{10-x}{2}$. Our ratio now becomes $\frac{8}{6} = \frac{x+ \frac{10-x}{2}}{DF}$. Solving for $DF$ gives us $DF = \frac{30+3x}{8}$. Since $DF$ is a height of the triangle $ADE$, $FE^2 + DF^2 = x^2$, or $DF = \sqrt{x^2 - (\frac{10-x}{2})^2}$. Solving the equation $\frac{30+3x}{8} = \sqrt{x^2 - (\frac{10-x}{2})^2}$ gives us $x = \frac{250}{39}$, so our answer is $250+39 = \boxed{289}$.

 



\textbf{Solution 3 (Algebra w/ Law of Cosines)}

As in the diagram, let $DE = x$. Consider point $G$ on the diagram shown above. Our goal is to be able to perform Pythagorean Theorem on $DG, GC$, and $DC$. Let $GE = \frac{10-x}{2}$. Therefore, it is trivial to see that $GC^2 = \Big(x + \frac{10-x}{2}\Big)^2$ (leave everything squared so that it cancels nicely at the end). By Pythagorean Theorem on Triangle $DGE$, we know that $DG^2 = x^2 - \Big(\frac{10-x}{2}\Big)^2$. Finally, we apply Law of Cosines on Triangle $DBC$. We know that $\cos(\angle DBC) = \frac{3}{5}$. Therefore, we get that $DC^2 = (10-x)^2 + 12^2 - 2(12)(10-x)\frac{3}{5}$. We can now do our final calculation:
 \[DG^2 + GC^2 = DC^2 \implies x^2 - \Big(\frac{10-x}{2}\Big)^2 + \Big(x + \frac{10-x}{2}\Big)^2 = (10-x)^2 + 12^2 - 2(12)(10-x)\frac{3}{5}\]After some quick cleaning up, we get
 \[30x = \frac{72}{5} + 100 \implies x = \frac{250}{39}\]Therefore, our answer is $250+39=\boxed{289}$.

 



\textbf{Solution 4 (Coordinates)}

Let $B = (0, 0)$, $C = (12, 0)$, and $A = (6, 8)$. Then, let $x$ be in the interval $0<x<2$ and parametrically define $D$ and $E$ as $(6-3x, 8-4x)$ and $(12-3x, 4x)$ respectively. Note that $AD = 5x$, so $DE = 5x$. This means that
 \begin{align*} \sqrt{36+(8x-8)^2} &= 5x\\ 36+(8x-8)^2 &= 25x^2\\ 64x^2-128x+100 &= 25x^2\\ 39x^2-128x+100 &= 0\\ x &= \dfrac{128\pm\sqrt{16384-15600}}{78}\\ x &= \dfrac{100}{78}, 2\\ \end{align*}
However, since $2$ is extraneous by definition, $x=\dfrac{50}{39}\implies AD = \dfrac{250}{39}\implies\boxed{289}$ ~ mathwiz0803

 



\textbf{Solution 5 (Law of Cosines)}

As shown in the diagram, let $x$ denote $\overline{AD}$. Let us denote the foot of the altitude of $A$ to $\overline{BC}$ as $F$. Note that $\overline{AE}$ can be expressed as $10-x$ and $\triangle{ABF}$ is a $6-8-10$ triangle . Therefore, $\sin(\angle{BAF})=\frac{3}{5}$ and $\cos(\angle{BAF})=\frac{4}{5}$. Before we can proceed with the Law of Cosines, we must determine $\cos(\angle{BAC})=\cos(2\cdot \angle{BAF})=\cos^2(\angle{BAF})-\sin^2(\angle{BAF})=\frac{7}{25}$. Using LOC, we can write the following statement:
 \[(\overline{DE})^2=(\overline{AD})^2+\overline{AE}^2-2(\overline{AD})(\overline{AE})\cos(\angle{BAC})\implies\]
 \[x^2=x^2+(10-x)^2-2(x)(10-x)\left(\frac{7}{25}\right)\implies\]
 \[0=(10-x)^2-\frac{14x}{25}(10-x)\implies\]
 \[0=10-x-\frac{14x}{25}\implies\]
 \[10=\frac{39x}{25}\implies x=\frac{250}{39}\]Thus, the desired answer is $\boxed{289}$ ~ blitzkrieg21

 



\textbf{Solution 6}

In isosceles triangle, draw the altitude from $D$ onto $\overline{AD}$. Let the point of intersection be $X$. Clearly, $AE=10-AD$, and hence $AX=\frac{10-AD}{2}$.

 Now, we recognise that the perpendicular from $A$ onto $\overline{AD}$ gives us two $6$-$8$-$10$ triangles. So, we calculate $\sin \angle ABC=\frac{4}{5}$ and $\cos \angle ABC=\frac{3}{5}$
$\angle BAC = 180-2\cdot\angle ABC$. And hence,

 $\cos \angle BAC = \cos \angle (180-2\cdot\angle ABC) = -\cos (2\cdot\angle ABC) = \sin^2 \angle ABC - \cos^2 \angle ABC = 2\cos^2 \angle ABC - 1 = \frac{32}{25}-\frac{25}{25}=\frac{7}{25}$

 Inspecting $\triangle ADX$ gives us $\cos \angle BAC = \frac{\frac{10-x}{2}}{x} = \frac{10-x}{2x}$Solving the equation $\frac{10-x}{2x}=\frac{7}{25}$ gives $x= \frac{250}{39} \implies\boxed{289}$




\textbf{Solution 7 (Area into Similar Triangles)}

After calling $x=AD=DE=EC$ and $10-x=AE=BD$, we see we have length ratios in terms of $x$, which motivates area ratios. We look at the area of triangle $ADC$ in two ways in order to find $DG$ (perpendicular from $D$ to $AB$), and then use similar triangles to find $x$. 

 Using area ratios, $[ADC] = \frac{x}{10}\cdot[ABC] = \frac{x}{10} \cdot 48 = \frac{24x}{5}$. (To find the total area $[ABC] = 48$, drop the altitudes from $A$ to $BC$, and call the foot of the altitude $F$. By the 6-8-10 triangle, the height $AF$ is $8$ and the area of $ABC$ is $48$.)

 The second way of finding the area of triangle $ACD$ is $\frac{1}{2}bh$. The base is $AC=10$, and $DG$ is the height. Therefore, 

 \begin{align*}[ACD] = \frac{24x}{5} &= \frac{1}{2} \cdot 10 \cdot DG \\[5pt]\frac{24x}{25} &= DG\end{align*}
 

 Now we have $DG$ in terms of $x$, we use the similar triangles $GCD$ and $FAC$ and set up the proportion\begin{align*}\frac{DG}{CF} &= \frac{GC}{FA} \\[5pt]\frac{\frac{24x}{25}}{6} &= \frac{5+\frac{x}{2}}{8} \\[5pt]x &= \frac{250}{39}.\end{align*}
So, our answer is $\boxed{289}$.\textbf{-epiconan}

 



\textbf{Solution 8 (Easiest way- Coordinates without bash)}

Let $B=(0, 0)$, and $C=(12, 0)$. From there, we know that $A=(6, 8)$, so line $AB$ is $y=\frac{4}{3}x$. Hence, $D=(a, \frac{4}{3}a)$ for some $a$, and $BD=\frac{5}{3}a$ so $AD=10-\frac{5}{3}a$. Now, notice that by symmetry, $E=(6+a, 8-\frac{4}{3}a)$, so $ED^2=6^2+(8-\frac{8}{3}a)^2$. Because $AD=ED$, we now have $(10-\frac{5}{3})^2=6^2+(8-\frac{8}{3}a)^2$, which simplifies to $\frac{64}{9}a^2-\frac{128}{3}a+100=\frac{25}{9}a^2-\frac{100}{3}a+100$, so $\frac{39}{9}a=\frac{13}{3}a=\frac{28}{3}$, and $a=\frac{28}{13}$.It follows that $AD=10-\frac{5}{3}\times\frac{28}{13}=10-\frac{140}{39}=\frac{390-140}{39}=\frac{250}{39}$, and our answer is $250+39=\boxed{289}$.

 -Stormersyle

 



\textbf{Solution 9 one second accurate solve(1 variable equation)}

Doing law of cosines we know that $\cos A$ is $\frac{7}{25}.$* Dropping the perpendicular from $D$ to $AE$ we get that 
 \[\frac{10-x}{2}=\frac{7x}{25}.\] Solving for $x$ we get $\frac{250}{39}$ so our answer is $289$.

 -harsha12345

 
            \begin{itemize}
        \item It is good to remember that doubling the smallest angle of a 3-4-5 triangle gives the larger (not right) angle in a 7-24-25 triangle.
            \end{itemize}
            



\textbf{Solution 10 (Law of Sines)}

Let's label $\angle A = \theta$ and $\angle ECD = \alpha$. Using isosceles triangle properties and the triangle angle sum equation, we get 
 \[180-(180-2\theta+\alpha) + \frac{180-\theta}{2} + \left(\frac{180-\theta}{2} - \alpha\right) = 180.\] Solving, we find $\theta = 2 \alpha$. 

 
Relabelling our triangle, we get $\angle ABC = 90 - \alpha$. Dropping an altitude from $A$ to $BC$ and using the Pythagorean theorem, we find $[ABC] = 48$. Using the sine area formula, we see $\frac12 \cdot 10 \cdot 12 \cdot \sin(90-\alpha) = 48$. Plugging in our sine angle cofunction identity, $\sin(90-\alpha) = \cos(\alpha)$, we get $\alpha = \cos{^{-1}}{\frac45}$. 

 
Now, using the Law of Sines on $\triangle ADE$, we get 
 \[\frac{\sin{2\alpha}}{\frac{p}{q}} = \frac{\sin{(180-4\alpha)}}{10-\frac{p}{q}}.\] After applying numerous trigonometric and algebraic tricks, identities, and simplifications, such as $\sin{(180-4\alpha)}=\sin{4\alpha}$ and $\sin{\left(\cos{^{-1}}{\frac45}\right)} = \frac35$, we find $\frac{p}{q} = \frac{10\sin{2\alpha}}{\sin{4\alpha}+\sin{2\alpha}} = \frac{250}{39}$. 

 


 Therefore, our answer is $250 + 39 = \boxed{289}$.

 



\textbf{Solution 11 (Trigonometry)}

We start by labelling a few angles (all of them in degrees). Let $\angle{BAC}=2\alpha = \angle{AED}, \angle{EDC}=\angle{ECD}=\alpha, \angle{DEC}=180-2\alpha, \angle{BDC}=3\alpha, \angle{DCB}=90-2\alpha, \angle{DBC}=90-\alpha$. Also let $AD=a$. By sine rule in $\triangle{ADE},$ we get $\frac{a}{\sin{2\alpha}}=\frac{10-a}{\sin{4\alpha}} \implies \cos{2\alpha}=\frac{5}{a}-\frac{1}{2}$Using sine rule in $\triangle{ABC}$, we get $\sin{\alpha}=\frac{3}{5}$. Hence we get $\cos{2\alpha}=1-2\sin^2{\alpha}=1-\frac{18}{25}=\frac{7}{25}$. Hence $\frac{5}{a}=\frac{1}{2}+\frac{7}{25}=\frac{39}{50} \implies a=\frac{250}{39}$. Therefore, our answer is $\boxed{289}$

 Alternatively, use sine rule in $\triangle{BDC}$. (It’s easier)

 



\textbf{Solution 12 (Double Angle Identity)}

We let $AD=x$. Then, angle $A$ is $2\sin^{-1}(\frac{3}{5})$ and so is angle $AED$. We note that $AE=10-x$. We drop an altitude from $D$ to $AE$, and we call the foot $F$. We note that $AF=\frac{10-x}{2}$. Using the double angle identity, we have $\sin(2\sin^{-1}(\frac{3}{5}))=2(\frac{3}{5})(\frac{4}{5})=\frac{24}{25}.$ Therefore, $DG=\frac{24}{25}AD.$ We now use the Pythagorean Theorem, which gives $(\frac{10-x}{2})^2+(\frac{24}{25}x)^2=x^2$. Rearranging and simplifying, this becomes $429x^2-12500x+62500=0$. Using the quadratic formula, this is $\frac{12500\pm\sqrt{12500^2-250000\cdot429}}{858}$. We take out a $10000$ from the square root and make it a $100$ outside of the square root to make it simpler. We end up with $\frac{12500\pm7000}{858}$. We note that this must be less than 10 to ensure that $10-x$ is positive. Therefore, we take the minus, and we get $\frac{5500}{858}=\frac{250}{39} \implies \fbox{289}.$



	\item (\textit{AIME I 2018}) For each ordered pair of real numbers $(x,y)$ satisfying

 \[\log_2(2x+y) = \log_4(x^2+xy+7y^2)\]there is a real number $K$ such that

 \[\log_3(3x+y) = \log_9(3x^2+4xy+Ky^2).\]Find the product of all possible values of $K$.


 
 \textbf{Answer: }189



\textbf{Solution 1}

Using the logarithmic property $\log_{a^n}b^n = \log_{a}b$, we note that 
 \[(2x+y)^2 = x^2+xy+7y^2\]That gives 
 \[x^2+xy-2y^2=0\] upon simplification and division by $3$. Factoring $x^2+xy-2y^2=0$ gives 
 \[(x+2y)(x-y)=0\] Then, 
 \[x=y \text{ or }x=-2y\]From the second equation, 
 \[9x^2+6xy+y^2=3x^2+4xy+Ky^2\] If we take $x=y$, we see that $K=9$. If we take $x=-2y$, we see that $K=21$. The product is $\boxed{189}$.

 -expiLnCalc

 



\textbf{Solution 2}

Do as done in Solution 1 to get 
 \[x^2+xy-2y^2=0\] 
 \[\implies (\frac{x}{y})^2+\frac{x}{y}-2=0\] 
 \[\implies \frac{x}{y}=\frac{-1\pm\sqrt{1+8}}{2}=1,-2\] Do as done in Solution 1 to get 
 \[9x^2+6xy+y^2=3x^2+4xy+Ky^2\] 
 \[\implies 6x^2+2xy+(1-K)y^2=0\] 
 \[\implies 6(\frac{x}{y})^2+2\frac{x}{y}+(1-K)=0\] 
 \[\implies \frac{x}{y}=\frac{-2\pm \sqrt{4-24(1-K)}}{12}\] 
 \[\implies \frac{x}{y}=\frac{-2\pm 2\sqrt{6K-5}}{12}=\frac{-1\pm \sqrt{6K-5}}{6}\]If $\frac{x}{y}=1$ then 
 \[1=\frac{-1\pm \sqrt{6K-5}}{6}\] 
 \[\implies 6=-1\pm \sqrt{6K-5}\] 
 \[\implies 7=\pm \sqrt{6K-5}\] 
 \[\implies 49=6K-5\] 
 \[\implies K=9\]If $\frac{x}{y}=-2$ then 
 \[-2=\frac{-1\pm \sqrt{6K-5}}{6}\] 
 \[\implies -12=-1\pm \sqrt{6K-5}\] 
 \[\implies -11=\sqrt{6K-5}\] 
 \[\implies 121=6K-5\] 
 \[\implies 126=6K\] 
 \[\implies K=21\]Hence our final answer is $21\cdot 9=\boxed{189}$-vsamc$\newline$




\textbf{Solution 3 (Official MAA)}

Because $x^2+xy+7y^2=\left(x+\frac{y}{2}\right)^2+\frac{27}{4}y^2>0,$ the right side of the first equation is real. It follows that the left side of the equation is also real, so $2x+y>0$ and 
 \[\log_2(2x+y)=\log_{2^2}(2x+y)^2=\log_4(4x^2+4xy+y^2).\] Thus $4x^2+4xy+y^2=x^2+xy+7y^2,$ which implies that $0=x^2+xy-2y^2=(x+2y)(x-y).$ Therefore either $x=-2y$ or $x=y,$ and because $2x+y>0,$ $x$ must be positive and $3x+y=x+(2x+y)>0.$ Similarly, 
 \[\log_3(3x+y)=\log_{3^2}(3x+y)^2=\log_9(9x^2+6xy+y^2).\] If $x=-2y\ne 0,$ then $9x^2+6xy+y^2=36y^2-12y^2+y^2=25y^2=3x^2+4xy+Ky^2$ when $K=21.$ If $x=y\ne 0,$ then $9x^2+6xy+y^2=16y^2=3x^2+4xy+Ky^2$ when $K=9.$ The requested product is $21\cdot9=189.$



	\item (\textit{AIME I 2018}) Let $N$ be the number of complex numbers $z$ with the properties that $|z|=1$ and $z^{6!}-z^{5!}$ is a real number. Find the remainder when $N$ is divided by $1000$.


 
 \textbf{Answer: }440



\textbf{Solution 1}

Let $a=z^{120}$. This simplifies the problem constraint to $a^6-a \in \mathbb{R}$. This is true if $\text{Im}(a^6)=\text{Im}(a)$. Let $\theta$ be the angle $a$ makes with the positive x-axis. Note that there is exactly one $a$ for each angle $0\le\theta<2\pi$. We are given $\sin\theta = \sin{6\theta}$. Note that $\sin \theta = \sin (\pi - \theta)$ and $\sin \theta = \sin (\theta + 2\pi)$. We can use these facts to create two types of solutions: 

 
 \[\sin \theta = \sin ((2m + 1)\pi - \theta)\]

 which implies that $(2m+1)\pi-\theta = 6\theta$ and reduces to $\frac{(2m + 1)\pi}{7} = \theta$. There are 7 solutions for this.

 
 \[\sin \theta = \sin (2n\pi + \theta)\]

 which implies that $2n\pi+\theta=6\theta$ and reduces to $\frac{2n\pi}{5} = \theta$. There are 5 solutions for this, totaling 12 values of $a$.

 For each of these solutions for $a$, there are necessarily $120$ solutions for $z$. Thus, there are $12\cdot 120=1440$ solutions for $z$, yielding an answer of $\boxed{440}$.

 



\textbf{Solution 2}

The constraint mentioned in the problem is equivalent to the requirement that the imaginary part is equal to $0$. Since $|z|=1$, let $z=\cos \theta + i\sin \theta$, then we can write the imaginary part of $\Im(z^{6!}-z^{5!})=\Im(z^{720}-z^{120})=\sin\left(720\theta\right)-\sin\left(120\theta\right)=0$. Using the sum-to-product formula, we get $\sin\left(720\theta\right)-\sin\left(120\theta\right)=2\cos\left(\frac{720\theta+120\theta}{2}\right)\sin\left(\frac{720\theta-120\theta}{2}\right)=2\cos\left(\frac{840\theta}{2}\right)\sin\left(\frac{600\theta}{2}\right)\implies \cos\left(\frac{840\theta}{2}\right)=0$ or $\sin\left(\frac{600\theta}{2}\right)=0$. The former yields $840$ solutions, and the latter yields $600$ solutions, giving a total of $840+600=1440$ solution, so our answer is $\boxed{440}$.

 



\textbf{Solution 3}

As mentioned in solution one, for the difference of two complex numbers to be real, their imaginary parts must be equal. We use exponential form of complex numbers. Let $z = e^{i \theta}$. We have two cases to consider. Either $z^{6!} = z^{5!}$, or $z^{6!}$ and $z^{5!}$ are reflections across the imaginary axis.If $z^{6!} = z^{5!}$, then $e^{6! \theta i} = e^{5! \theta i}$. Thus, $720 \theta = 120 \theta$ or $600\theta = 0$, giving us 600 solutions. (Equalities are taken modulo $2 \pi$)For the second case, $e^{6! \theta i} = e^{(\pi - 5!\theta)i}$. This means $840 \theta = \pi$, giving us 840 solutions.Our total count is thus $1440$, yielding a final answer of $\boxed{440}$.

 



\textbf{Solution 4}

Because $|z| = 1,$ we know that $z\overline{z} = 1^2 = 1.$ Hence $\overline{z} = \frac 1 {z}.$ Because $z^{6!}-z^{5!}$ is real, it is equal to its complex conjugate. Hence $z^{6!}-z^{5!} = \overline{z^{6!}}-\overline{z^{5!}}.$ Substituting the expression we that we derived earlier, we get $z^{720}-z^{120} = \frac 1{z^{720}} - \frac 1{z^{120}}.$ This leaves us with a polynomial whose leading term is $z^{1440}.$ Hence our answer is $\boxed{440}$.~Shen Kislay Kai 2022

 Note: This is actually not rigorous, because how to we know that all of the roots of such a polynomial are distinct? One can proceed as follows. Factoring gives us that $(z^{840}+1)(z^{600}-1)=0,$ so this implies that $z^{840} = -1$ OR $z^{600}=1.$ To show no $z$ satisfies both of these conditions, notice that if $w^{840} = -1$ and $w^{600}=1$ for some complex number $w$, then $w^{240} = w^{840-600} = -1,$ which implies that $w^{360} = w^{600-240}=-1,$ which implies that $w^{120}=w^{360-240}=1.$ This is a contradiction since then $w^{240}$ would also have to equal $(w^{120})^2 = 1.$ Therefore the total number of solutions is $840+600 = 1\boxed{440}.$Minor Edit by ~ hazio Jheng

 I'm pretty sure it works since fundamental theory of algebra counts multiplicity.Above, even though there are 1440 solutions if some solutions are the same they should not be counted twice in the working cases for z

 



\textbf{Solution 5}

Since $|z|=1$, let $z=\cos \theta + i\sin \theta$. For $z^{6!}-z^{5!}$ to be real, the imaginary parts of $z^{6!}$ and $z^{5!}$ must be equal, so $\sin 720\theta=\sin 120\theta$. We need to find all solutions for $\theta$ in the interval $[0,2\pi)$. This can be done by graphing $y=\sin 720\theta$ and $y=\sin 120\theta$ and finding their intersections. Since the period of $y=\sin 720\theta$ is $\frac{\pi}{360}$ and the period of $y=\sin 120\theta$ is $\frac{\pi}{60}$, the common period of both graphs is $\frac{\pi}{60}$. Therefore, we only graph the functions in the domain $[0, \frac{\pi}{60})$. We can clearly see that there are twelve points of intersection. However, since we only graphed $\frac{1}{120}$ of the interval $[0,2\pi)$, we need to multiply our answer by $120$. The answer is $12 \cdot 120 = 1440 = \boxed{440} \pmod{1000}$.

 



\textbf{Solution 6 (Official MAA)}

If $z$ satisfies the given conditions, there is a $\theta \in [0,2\pi]$ such that $z=e^{i\theta}$ and $e^{720\theta i-120\theta i}$ is real. This difference is real if and only if either the two numbers $720\theta$ and $120\theta$ represent the same angle or the two numbers represent supplementary angles. In the first case there is an integer $k$ such that $720\theta=120\theta+2k\pi,$ which implies that $\theta$ is a multiple of $\frac{\pi}{300}.$ In the second case there is an integer $k$ such that $720\theta=-120\theta+(2k+1)\pi,$ which implies that $\theta$ is $\frac{\pi}{840}$ plus a multiple of $\frac{\pi}{420}.$ In the interval $[0,2\pi]$ there are $600$ values of $\theta$ that are multiples of $\frac{\pi}{300},$ there are $840$ values that are $\frac{\pi}{840}$ plus a multiple of $\frac{\pi}{420},$ and there are no values of $\theta$ that satisfy both of these conditions. Therefore theremust be $600+840=1440$ complex numbers satisfying the given conditions. The requested remainder is $440.$




\textbf{Solution 7}

$z^{720}-z^{120}=z^{120}(z^{600}-1)$ firstly we consider that $z^{600}-1=0$ obviously there are $600$ different roots

 Then, we consider that $z^{720}=a+bi, z^{120}=-a+bi$ or $z^{720}=-a-bi, z^{120}=a-bi$ which leads to $840\theta=\pi+2k\pi, k\in N$ so there are $840$ roots, in all, $840+600 \equiv 440 ($ mod $1000)$.

 


	\item (\textit{AIME I 2018}) A right hexagonal prism has height $2$. The bases are regular hexagons with side length $1$. Any $3$ of the $12$ vertices determine a triangle. Find the number of these triangles that are isosceles (including equilateral triangles).


 
 \textbf{Answer: }052



\textbf{Solution 1}

We can consider two cases: when the three vertices are on one base, and when the vertices are on two bases.

 Case 1: vertices are on one base. Then we can call one of the vertices $A$ for distinction. Either the triangle can have sides $1, 1, \sqrt{3}$ with 6 cases or $\sqrt{3}, \sqrt{3}, \sqrt{3}$ with 2 cases. This can be repeated on the other base for $16$ cases.

 Case 2: The vertices span two bases. WLOG call the only vertex on one of the bases $X$. Call the closest vertex on the other base $B$, and label clockwise $C, D, E, F, G$. We will multiply the following scenarios by $12$, because the top vertex can have $6$ positions and the top vertex can be on the other base. We can have $XCG, XDF$, but we are not done! Don't forget that the problem statement implies that the longest diagonal in a base is $2$ and the height is $2$, so $XBE$ is also correct! Those are the only three cases, so there are $12*3=36$ cases for this case.

 In total there's $\boxed{052}$ cases.

 



\textbf{Solution 2}

If there are two edges on a single diameter, there would be six diameters. There are four ways to put the third number, and four equilateral triangles. There are $4+ 2 \cdot 2\cdot 6 = 28$ ways. Then, if one length was $\sqrt{3}$ but no side on the diameter, there would be twelve was to put the $\sqrt{3}$ side, and two ways to put the other point. $2 \cdot 12 = 24$ for four ways to put the third point. Adding the number up, the final answer is $24+28 = \boxed{052}.$




\textbf{Solution 3}

To start, let's find the distances from any vertex (call it A, doesn't matter since the prism is symmetrical) to all the other 11 vertices. Using Pythagorean relations, we find that the distances are $1$, $\sqrt{3}$, $2$, $\sqrt{3}$, $1$, $\sqrt{7}$, $\sqrt{5}$, $2$, $\sqrt{5}$, $\sqrt{7}$, and $\sqrt{8}$.

 We can clearly form an isosceles triangle using any vertex and any two neighboring edges that are equal. There are 12 total vertices, all of which are symmetrical to the vertex A. Thus, for each vertex, we can form 5 isosceles triangles (with neighboring edges $1$-$1$, $2$-$2$, $\sqrt{3}$-$\sqrt{3}$, $\sqrt{5}$-$\sqrt{5}$, and $\sqrt{7}$-$\sqrt{7}$) , so have so far $12 \times 5 = 60$ isosceles triangles.

 To check for overcounting, note that isosceles triangles can be split into two major categories: equilateral isosceles, and non-equilateral isosceles. We only counted non-equilateral isosceles triangles once since there is only one vertex whose two neighboring edges are equal. But equilateral triangles were counted three times (namely once for each vertex).

 By observation, there are only 4 equilateral triangles (1-1-1 side lengths), two on each hexagonal face. Since we counted each two more times than we should of (three instead of once), we will subtract ($4$ eq. triangles) $\times$ ($2$ overcounts per eq. triangle) = $8$ overcounts.

 Finally, we have $60 - 8$ overcounts = $\boxed{052}$ isoceles triangles.

 



\textbf{Solution 4}

Start by drawing the hexagonal prism. Start from one of the points, let's call it $A$. From simple inspection, we can see that each point we choose creates $4$ isosceles triangles; $A$ and the two adjacent points, $A$ and the two points adjacent to those, $A$ and the adjacent point on the other face, and $A$ and the two points adjacent to those. Here, we have $12\cdot4\Rightarrow48$ isosceles triangles.

 Next, notice that you can create an isosceles triangle with the bases of the hexagonal prism. Let's say that our base, starting from the very left-most point of the hexagon, has points $A, B, C, D, E, F$, rotating counter-clockwise. We can create a triangle with either points $\bigtriangleup{BCE}$ or $\bigtriangleup{BCF}$. WLOG, assuming we use $\bigtriangleup{BCE}$, we can use the Law of Cosines and find that $\overline{CE}$ = $\sqrt{3}$. Thus, this would mean $\overline{CF}$ is $2$, which is the exactly the same length as the height of the hexagonal prism. There are $4$ possibilities of this case, so our final solution is $48+4=\boxed{052}$.

 Solution by IronicNinja~

 



\textbf{Solution 5 (Official MAA)}

There are $2$ equilateral and $6$ other isosceles triangles on each base, providing a total of $16$ triangles when all $3$ vertices are chosen from the same base. To count the rest of the isosceles triangles, there are $2$ ways to choose the base that will contain $2$ vertices. Assume that $2$ vertices have been chosen from the bottom base. Then there are no isosceles triangles if the $2$ vertices on the bottom base are adjacent vertices of the hexagon, $6\cdot 2=12$ isosceles triangles if the $2$ vertices on the bottom base have $1$ vertex between them on the hexagon, and $3\cdot2=6$ isosceles triangles if the $2$ vertices on the bottom base have $2$ vertices between them on the hexagon.

 The total number of isosceles triangles is $16+2(12+6)=52.$




\textbf{Solution 6 (easy)}

possible options for the side lengths of the triange are:$(2,2,\sqrt8); (\sqrt3, \sqrt3, \sqrt3); (\sqrt3, \sqrt5, \sqrt5); (\sqrt3, \sqrt7, \sqrt7); (1,1,\sqrt3).$do casework to get $12+12+12+12+4$ where $4$ is for the number of equilateral triangles.

 


	\item (\textit{AIME I 2018}) Let $ABCDEF$ be an equiangular hexagon such that $AB=6, BC=8, CD=10$, and $DE=12$. Denote $d$ the diameter of the largest circle that fits inside the hexagon. Find $d^2$.


 
 \textbf{Answer: }147



\textbf{Solution 1}


                    \begin{figure}[H]
                        \centering
                        \includegraphics[width=0.5\textwidth]{Solution/AIME_I_2018_8_Solution_Figure_1.png}
                    \end{figure}
                    First of all, draw a good diagram! This is always the key to solving any geometry problem. Once you draw it, realize that $EF=2, FA=16$. Why? Because since the hexagon is equiangular, we can put an equilateral triangle around it, with side length $6+8+10=24$. Then, if you drew it to scale, notice that the "widest" this circle can be according to $AF, CD$ is $7\sqrt{3}$.  And it will be obvious that the sides won't be inside the circle, so our answer is $\boxed{147}$.

 -expiLnCalc

 



\textbf{Solution 2}

Like solution 1, draw out the large equilateral triangle with side length $24$. Let the tangent point of the circle at $\overline{CD}$ be G and the tangent point of the circle at $\overline{AF}$ be H. Clearly, GH is the diameter of our circle, and is also perpendicular to $\overline{CD}$ and $\overline{AF}$.

 The equilateral triangle of side length $10$ is similar to our large equilateral triangle of $24$. And the height of the former equilateral triangle is $\sqrt{10^2-5^2}=5\sqrt{3}$. By our similarity condition,$\frac{10}{24}=\frac{5\sqrt{3}}{d+5\sqrt{3}}$

 Solving this equation gives $d=7\sqrt{3}$, and $d^2=\boxed{147}$



	\item (\textit{AIME I 2018}) Find the number of four-element subsets of $\{1,2,3,4,\dots, 20\}$ with the property that two distinct elements of a subset have a sum of $16$, and two distinct elements of a subset have a sum of $24$. For example, $\{3,5,13,19\}$ and $\{6,10,20,18\}$ are two such subsets.


 
 \textbf{Answer: }210



\textbf{Solution 1}

This problem is tricky because it is the capital of a few "bashy" calculations. Nevertheless, the process is straightforward. Call the set $\{a, b, c, d\}$.

 Note that there are only two cases:

 $\textrm{Case } 1 \textrm{:}$ where $a + b = 16$ and $c + d = 24$.

 $\textrm{Case } 2 \textrm{:}$ where $a + b = 16$ and $a + c = 24$.

 Also note that there is no overlap between the two situations! This is because if they overlapped, adding the two equations of both cases and canceling out gives you $a=d$, which cannot be true.

 $\textrm{Case } 1 \textrm{:}$ This is probably the simplest: just make a list of possible combinations for $\{a, b\}$ and $\{c, d\}$. We get $\{1, 15\}\dots\{7, 9\}$ for the first and $\{4, 20\}\dots\{11, 13\}$ for the second. That appears to give us $7*8=56$ solutions, right? No, as because elements cannot repeat, take out the supposed sets:

 
 \[\begin{array}{l} \{1,15,9,15\}, \{2,14,10,14\}, \{3,13,11,13\}, \{4,12,4,20\}, \{5,11,5,19\},\\ \{5,11,11,13\}, \{6,10,6,18\}, \{6,10,10,14\}, \{7,9,9,15\}, \{7,9,7,17\}. \end{array}\]

 Thus, $10$ cases can be removed. So $46$ for $\textrm{Case } 1$.

 $\textrm{Case } 2 \textrm{:}$ We can look for solutions by listing possible $a$ values and filling in the blanks. Start with $a=4$, as that is the minimum. We find $\{4, 12, 20, ?\}$, and likewise up to $a=15$. But we can't have $a=8$ or $a=12$ because $a=b$ or $a=c$, respectively! Now, it would seem like there are $10$ values for $a$ and $17$ unique values for each $?$, giving a total of $170$, but that is once again not true because there are some repeated values!There are two cases of overcounting:

 $\textrm{Case } a$: Some sets are counted twice because of the ordering of $a$ and $d$.

 $(5,11,13,19)$ and $(5,11,19,13)$

 The same is true for $(6,10,14,18)$ and $(7,9,15,17)$.

 $\textrm{Case } b$: Some sets are counted twice because they have the same $b$ and $c$ values.

 This case includes:

 $(1,15,9,7)$ and $(7,9,15,1)$
$(2,14,10,6)$ and $(6,10,14,2)$
$(3,13,11,5)$ and $(5,11,13,3)$

 So we need to subtract $6$ overcounts. Thus, there are $170-6=164$ solutions for $\textrm{Case } 2$.

 Total gives $\boxed{210}$.

 -expiLnCalc

 -Builtin

 



\textbf{Solution 2}

Let's say our four elements in our subset are $a,b,c,d$. We have two cases. Note that the order of the elements / the element letters themselves don't matter since they are all on equal grounds at the start.

 
$\textrm{Case } 1 \textrm{:}$ $a+b = 16$ and $c+d = 24$. 

 List out possibilities for $a+b$ $(\text{i.e. } 1+15, 2+14, 3+13 \text{ etc.})$ but don't list $8+8$ because those are the same elements and that is restricted. 

 Then list out the possibilities for $c+d \text{ }(\text{i.e. } 4+20, 5+19, 6+18, \text{ etc.})$ but don't list $12+12$ because they are the same elements. 

 This will give you $7 \cdot 8$ elements, which is $56$. However, as stated above, we have overlap. Just count starting from $a+b$.  $15,14,13,11,10,9,7,6,5,4$ all overlap once, which is $10$, thus $56 - 10 = 46$ cases in this case. Do not include $8$ or $12$ because if they are included, then two elements will be equal.

 There cannot be any overlap in this case. For there to be an overlap let us choose the elements $a$ and $c$ (one from equation) without loss of generality (abbreviated WLOG). Then, let their sum be $16$ (one sum) WLOG. $a+c=16$ cannot also occur because this leads to $b=c$.

 


 $\textrm{Case } 2 \textrm{:}$ $a+b = 16$ and $b+c = 24$. 

 Here, $b$ is included in both equations. We can easily see that $a, b, c$ will never equal each other. 

 Furthermore, there are 17 choices for $d$ ($20 - 3$ included elements) for each $b$. Listing out the possible $b$s, we go from $15,14,13,11,10,9,7,6,5,4$. Do not include $8$ or $12$ because if they are included, then $a/c$ will be the same as $b$, which is restricted. 

 There are $10$ options there, and thus $10 \cdot 17 = 170$. But, note that if $d = b+8$, $a+d = a+b+8 = 24$, and so we have a double-counted set. Starting with $b=15$, we have $15, 14, 13, 11, 10, 9$ (where $d$ is $7, 6, 5, 3, 2, 1)$. That means there are $6$ double-counted cases. Thus $170 - 6 = 164$ cases in this case.

 Adding these up, we get $46+164 = \boxed{210}.$




\textbf{Solution 3 (Official MAA)}

There are two types of $\{a,b,c,d\} \subseteq \{1,2,3,4\dots,20\}$ that have the needed property. There is either an assignment of distinct values for $a,\,b,\,c,$ and $d$ such that $a+b=16$ and $c+d=24$ or an assignment such that $a+b=16$ and $a+c=24.$ These two types are mutually exclusive because $c+d=24$ and $a+c=24$ imply that $a=d.$ For the first type, there are $7$ choices for $\{a,b\},$ namely $\{1,15\},\,\{2,14\},\,\{3,13\},\,\{4,12\},\,\{5,11\},\,\{6,10\},$ and $\{7,9\},$ and there are $8$ choices for $\{c,d\},$ namely $\{4,20\},\,\{5,19\},\,\{6,18\},\,\{7,17\},\,\{8,16\},\,\{9,15\},\{10,14\},$ and $\{11,13\}.$ Thus a four-element subset of the first type can be formed by taking the union of one of $7$ two-element subsets with one of $8$ two-element subsets as long as these two subsets are disjoint. There are $10$ such pairings that are not disjoint out of the $7\cdot 8=56$ pairings, so there are $56-10=46$ subsets of the first type.

 For subsets of the second type, there are $10$ choices for the value of $a$ $(4,\,5,\,6,\,7,\,9,\,10,\,11,\,13,\,14,\,15)$ such that $b=16-a$ and $c=24-a$ can be two other elements of the subset. Note that in each of these cases, $c-b=(24-a)-(16-a)=8.$ For each of these, there are $20-3=17$ other values that can be chosen for the element $d$ in the subset. But $10\cdot 17=170$ counts some subsets more than once. In particular, a subset is counted twice if $b+d=24$ or $c+d=16$. In such cases either $d=a+8$ or $d=a-8$. There are exactly $6$ subsets where the role of $a$ can be played by two different elements of the set. They are $\{1,7,9,15\},\,\{2,6,10,14\},\,\{3,5,11,13\},\,\{5,11,13,19\},\,\{6,10,14,18\},$ and $\{7,9,15,17\}$. Thus there are $170-6=164$ subsets of the second type.

 In all, there are $46+164=210$ subsets with the required property.

 


	\item (\textit{AIME I 2018}) The wheel shown below consists of two circles and five spokes, with a label at each point where a spoke meets a circle. A bug walks along the wheel, starting at point $A$. At every step of the process, the bug walks from one labeled point to an adjacent labeled point. Along the inner circle the bug only walks in a counterclockwise direction, and along the outer circle the bug only walks in a clockwise direction. For example, the bug could travel along the path $AJABCHCHIJA$, which has $10$ steps. Let $n$ be the number of paths with $15$ steps that begin and end at point $A.$ Find the remainder when $n$ is divided by $1000$.


 
                    \begin{figure}[H]
                        \centering
                        \includegraphics[width=0.5\textwidth]{Problem/AIME_I_2018_10_Problem_Figure_1.png}
                    \end{figure}
                    

 \textbf{Answer: }004



\textbf{Solution 1}

We divide this up into casework.  The "directions" the bug can go are $\text{Clockwise}$, $\text{Counter-Clockwise}$, and $\text{Switching}$.  Let an $I$ signal going clockwise (because it has to be in the \textit{inner} circle), an $O$ signal going counter-clockwise, and an $S$ switching between inner and outer circles.  An example string of length fifteen that gets the bug back to $A$ would be $ISSIIISOOSISSII$.For the bug to end up back at $A$, the difference between the number of $I$'s and $O$'s must be a multiple of $5$.

 
            \begin{enumerate}[label=(\roman*)]
                <dt>Case 1  --  There are 15 more $I$'s than $O$'s.</dt>\item There is clearly $1$ way for this to happen.
            \end{enumerate}
            
            \begin{enumerate}[label=(\roman*)]
                <dt>Case 2  --  There are $5$ more $I$'s than $O$'s.</dt><dd>We split this case up into several sub-cases based on the number of $S$'s.
            \begin{enumerate}[label=(\roman*)]
                <dt>Sub-case 1  --  There are $10$ $S$'s and $5$ $I$'s.</dt>\item Notice that the number of ways to order the $I$'s and $O$'s are independent assortments because the $I$'s must be in the "even" spaces between $S$'s (i.e. before the 1st $S$, between the 2nd and 3rd $S$'s, etc.), while the $O$'s must be in the "odd" spaces.
            \end{enumerate}
            </dd>
            \end{enumerate}
            
            \begin{enumerate}[label=(\roman*)]
                <dd>
            \begin{enumerate}[label=(\roman*)]
                \item There are $6$ places to put the $I$'s (after the 0th, 2nd, 4th, 6th, 8th, and 10th $S$'s), and $4$ places to put the (0) $O$'s.  We use stars and bars to get an answer of $\binom{10}{5}\binom{4}{0}$<dt>Sub-case 2  --  There are $8$ $S$'s, $6$ $I$'s, and $1$ $O$.</dt>\item Similarly and by using stars and bars, we get an amount of $\binom{10}{4}\binom{4}{1}$
            \end{enumerate}
            </dd>\item All the other  sub-cases are similar, with a total of $\binom{10}{5}\binom{4}{0}+\binom{10}{4}\binom{4}{1}+\cdots+\binom{10}{1}\binom{4}{4}=\binom{14}{5}=2002$ by \href{https://artofproblemsolving.com/wiki/index.php?title=Combinatorial\_identity\#Vandermonde.27s\_Identity}{Vandermonde's Identity}.
            \end{enumerate}
            
            \begin{enumerate}[label=(\roman*)]
                <dt>Case 3  --  There are $5$ more $O$'s than $I$'s.</dt>\item This case is similar to the other case.<dd>Here is an example of a sub-case for this case.
            \begin{enumerate}[label=(\roman*)]
                <dt>Sub-case</dt>\item There are $10$ $S$'s and $5$ $O$'s.\item There are $\binom{9}{4}\binom{5}{0}$ ways to do this.
            \end{enumerate}
            </dd>\item We can see now that the pattern is going to be $\binom{9}{4}\binom{5}{0}+\binom{9}{3}\binom{5}{1}+\cdots+\binom{9}{0}\binom{5}{4}=\binom{14}{4}=1001$.
            \end{enumerate}
            
So, the total number of ways is $1+2002+1001=3004$ which gives $\boxed{004}$ as the answer.

 



\textbf{Solution 2 (Official MAA)}

Note that the set of sequences of moves the bug makes is in bijective correspondence with the set of strings of $X$s and $Y$s of length $15$, where $X$ denotes a move which is either counterclockwise or inward along a spoke and $Y$ denotes a move which is either clockwise or outward along a spoke.  (The proof of this basically boils down to the fact that which one depends on whether the bug is on the inner wheel or the outer wheel.)  Now the condition that the bug ends at A implies that the difference between the number of $X$s and the number of $Y$s is a multiple of $5$, and so we must have either $4$, $9$, or $14$ $X$s within the first fourteen moves with the last move being an $X$.  This implies the answer is 
 \[\binom{14}4+\binom{14}9+\binom{14}{14} = 3004\equiv \boxed{004}\pmod{1000}.\]




\textbf{Solution 3 (Similar To Solution 2 But Modified To Clarify)}

Let an $O$ signal a move that ends in the outer circle and $I$ signal a move that ends in the inner circle. Now notice that for a string of $15$ moves to end at $A$, the difference between $O$'s and $I$'s in the string must be a multiple of $5,$ and this is because the net displacement of the bug after its journey must be a multiple of $5$, as that is one cycle around the wheel.

 $15$ $I$'s: Trivially $1$ case.

 $5$ $O$'s and $10$ $I$'s: Since the string has to end in an $I$ for the bug to land on $A$, there are a total of $\binom{14}{5}=2002$ ways to put $5$ $O$'s in the first $14$ moves.

 $10$ $O$'s and $5$ $I$'s: Similarly there are $\binom{14}{4}=1001$ ways to put $5-1=4$ $I$'s in the first $14$ moves.

 $15$ $O$'s: Impossible since the string has to end with an I.

 This brings us an answer of $1+2002+1001=3004 \equiv \boxed{004} \pmod{1000}$.

 -Solution by mathleticguyyyyyyyyy-~Edited by ike.chen~Edited again by nevergonnagiveup

 



\textbf{Solution 4 (Recursion)}

Define $A_n$ to be the number of sequences of length $n$ that ends at $A$ and similarly for the other spokes. Also, let
 \[S_n=A_n+B_n+C_n+D_n+E_n+F_n+G_n+H_n+I_n+J_n\]At every point in the path, the bug has $2$ choices for its next move. Since $S_n$ represents the total number of paths we have $S_n=2^n$. We can create a recursive formula for $A_n:$

 \begin{align*} A_n&=J_{n-1}+E_{n-1} \\ &=(I_{n-2}+A_{n-2})+(D_{n-2}+F_{n-2}) \\ &=(H_{n-3}+B_{n-3})+(J_{n-3}+E_{n-3})+(G_{n-3}+C_{n-3})+(J_{n-3}+E_{n-3}) \\ &=(B_{n-3}+C_{n-3}+E_{n-3}+G_{n-3}+H_{n-3}+J_{n-3})+(J_{n-3}+E_{n-3}) \\ &=(S_{n-3}-(I_{n-3}+A_{n-3}+D_{n-3}+F_{n-3}))+A_{n-2} \\ &=2^{n-3}-(J_{n-2}+E_{n-2})+A_{n-2} \\ &=2^{n-3}-A_{n-1}+A_{n-2} \\ \end{align*}
$A_0=0$, $A_1=0$, $A_2=1$, $A_3=0$, $A_4=3$. Continuing the process yields $A_{15}=3\boxed{004}$.

 



\textbf{Solution 5 (Stars and Bars)}

The points on the outer circle and the points on the inner circle have a one-to-one correspondence. Therefore, it is equivalent to the scenario of switching directions back and forth in the inner circle but changing directions to make up for travelling on spokes. If the ant changes directions $s$ times, the number of moves the ant makes is $15-s$. Also note that is switches directions an even number of times, so we think of it in pairs:

 Case 1: 0 pairs

 The ant would have to only be in the inner circle and go in one direction. Hence there is $1$ way.

 Case 2: 1 pair

 This is equivalent to the ant making 13 moves in the inner circle where it switches directions 2 times (counterclockwise-clockwise-counterclockwise). Note that the number of counterclockwise and clockwise moves must be equivalent$\pmod{5}$. Let the number of counterclockwise moves be $C$ and clockwise moves be $C'$. The only two feasible solutions for $(C,C')$ are $(4,9)$ and $(9,4)$. Therefore, we use stars and bars to obtain $10+5=15$ ways for this case. 

 Case 3: 2 pairs

 The ant changes directions 4 times and makes 11 moves, so we obtain $(3,8)$ and $(8,3)$ using similar logic. Using stars and bars once again, we obtain $90+180=270$ ways for this case.

 Case 4: 3 pairs

 Again, the only solutions for $(C,C')$ are $(2,7)$ and $(7,2)$. Using stars and bars, we obtain $360+720=1080$.

 Case 5: 4 pairs

 The only solutions are $(1,6)$ and $(6,1)$. Using stars and bars, we obtain $840+420=1260$.

 Case 6: 5 pairs

 The only solutions are $(0,5)$ and $(5,0)$. Using stars and bars, we obtain $252+126=378$.Note that there are no more cases since $C$ and $C'$ have to be equivalent$\pmod{5}$, which would need to make $C$ or $C'$ negative. Adding all the cases, we get $1+15+270+1080+1260+378=3004$. Taking modulo $1000$, we obtain $3004 \equiv \boxed{004} \pmod{1000}$

 
The process of using stars and bars was the number of ways to group the counterclockwise and clockwise moves. For instance, if there are $p$ pairs of switches, then based on the solutions for $(C,C')$, we can divide $C$ into $p+1$ subsets and $C'$ into $p$ subsets. This is equivalent to putting $C$ objects in $p+1$ bins and $C'$ objects into $p$ bins where a bin can be empty which is modeled by consecutive switches. Then, we multiply both results. Therefore, for each $p$, we compute $\dbinom{C+p}{p}\dbinom{C'+p-1}{p-1}$ and add them up.

 



\textbf{Solution 6}

We can view the diagram as a collection of points on the complex plane that are connected by transformations/symmetries. The inner circle can be represented as the fifth roots of unity(WLOG let $A$ correspond to $1$) where moving between the inner circle can be mapped to multiplying by $e^{\frac{2i\pi}{5}}$ we can also view moving between circles as a multiplying by $-e^{-2i\pi} = -1.$ 

 Therefore the generating function for the paths beginning at $1$ or $A$ can be seen as $(e^{\frac{2i\pi}{5}} - 1)^{15}.$

 We want the positive real terms or the ones that end on $1$ and those add up to $\dbinom{15}{5}(e^{2i\pi/5})^5(-1)^{10} + \dbinom{15}{15}(e^{2i\pi/5})^{15} = 3004 \equiv \boxed{004} \pmod{1000}.$




\textbf{Solution 7 (Official MAA, clarified)}

Our motivation is that moving along the outer circle "counters" moving along the inner circle in terms of angular position. Hence, denote $I$ as moving along the inner circle OR moving from the outside circle to the inner circle. Also denote $O$ as moving along the outer circle OR moving from the inside circle to the outside circle.

 
Notice that eventually, moves that go from the inner circle to the outer circle must equal to moves that go from the outer circle from the inner circle. This is because we must always end up at the inner circle at the last move, meaning that every time we moved to the outer circle, we must have also eventually moved back to the inner circle. 

 
With this in mind, we can now "ignore" moving between inner and outer circle. We find that $O \equiv I \pmod 5$, since we must end up at our original angular position, and taking 5 moves in one direction lands us back in the same position.

 
Using this fact, we form $3$ cases:

 
$O - I = 5$, $O + I = 15$   which gives us $O = 10$. Since we have $10$ $O$s and $5$ $I$s, but the last move must be an $I$, we have $\binom{14}{10}=1001$ ways for this case.

 $O - I = -5$, $O + I = 15$   which gives us $O = 5$. We have $\binom{14}{5}=2002$ ways for this case.

 $O - I = -15$, $O + I = 15$   which gives us $O = 0$. (aka only moving in the inner circle). We have $\binom{14}{0}=1$ ways for this case.

 
Ultimately, the total number of ways is $1001 + 2002 + 1 = 3004$, and our answer is $\boxed{004}$.

 



\textbf{Solution 8}

Denote $I(N)$ to be the number of moves upon the inner ring, $O(N)$ to be defined analogously upon the outer ring. If truly we start at $A$ and end at $A$, then the number of moves on the exterior of the ring is a multiple of 5, analogously defined for the interior. The only exception is the paths from the exterior to the interior, which contribute 1. At maximum because periods of cyclicity 5 exist we require moves of multiples 5. Now if we were to consider an arrangement moving upon the outer and inner rings 0, 5, 10, 15 times each, then we have to append either 15, 10, 5, or 0 movements from the interior to the exterior. This is just stars and bars upon adding, which can be done in $\binom{14}{14} + \binom{14}{9} + \binom{14}{4} + \binom{14}{-1} = 1 + 2002 + 1001 = 3004$ plausible configurations. The number modulo 1000 is $\boxed{4}$.

 If one finds this hard to understand, we suppose that the combined number of movements upon $I(N) + O(N)$ contribute 0, 5, 10, or 15 to the total count as then we will always end up at the 0 modular 5 position, which is $A$. In order to accommodate for $I(N)$ we have to add back in the respective values of moving from $I(N)$ to $O(N)$. 

 


	\item (\textit{AIME I 2018}) Find the least positive integer $n$ such that when $3^n$ is written in base $143$, its two right-most digits in base $143$ are $01$.


 
 \textbf{Answer: }195



\textbf{Solution 1}

Note that the given condition is equivalent to $3^n \equiv 1 \pmod{143^2}$ and $143=11\cdot 13$. Because $\gcd(11^2, 13^2) = 1$, the desired condition is equivalent to $3^n \equiv 1 \pmod{121}$ and $3^n \equiv 1 \pmod{169}$.

 If $3^n \equiv 1 \pmod{121}$, one can see the sequence $1, 3, 9, 27, 81, 1, 3, 9...$ so $5|n$.

 Now if $3^n \equiv 1 \pmod{169}$, it is harder. But we do observe that $3^3 \equiv 1 \pmod{13}$, therefore $3^3 = 13a + 1$ for some integer $a$. So our goal is to find the first number $p_1$ such that $(13a+1)^ {p_1} \equiv 1 \pmod{169}$. Then, $p_1 \equiv 0 \pmod{13},$ which follows from the binomial theorem. It is not difficult to see that the smallest $p_1=13$, so ultimately $3^{39} \equiv 1 \pmod{169}$. Therefore, $39|n$.

 The first $n$ satisfying both criteria is thus $5\cdot 39=\boxed{195}$.

 -expiLnCalc

 



\textbf{Solution 2}

Note that Euler's Totient Theorem would not necessarily lead to the smallest $n$ and that in this case that $n$ is greater than $1000$. 

 We wish to find the least $n$ such that $3^n \equiv 1 \pmod{143^2}$. This factors as $143^2=11^{2}*13^{2}$. Because $\gcd(121, 169) = 1$, we can simply find the least  $n$ such that $3^n \equiv 1 \pmod{121}$ and $3^n \equiv 1 \pmod{169}$.

 Quick inspection yields $3^5 \equiv 1 \pmod{121}$ and $3^3 \equiv 1 \pmod{13}$. Now we must find the smallest $k$ such that $3^{3k} \equiv 1 \pmod{169}$. Euler's gives $3^{156} \equiv 1 \pmod{169}$. So $3k$ is a factor of $156$. This gives $k=1,2, 4, 13, 26, 52$. Some more inspection yields $k=13$ is the smallest valid $k$. So $3^5 \equiv 1 \pmod{121}$ and $3^{39} \equiv 1 \pmod{169}$. The least $n$ satisfying both is $lcm(5, 39)=\boxed{195}$. (RegularHexagon)

 



\textbf{Solution 3 (Big Bash)}

Listing out the powers of $3$, modulo $169$ and modulo $121$, we have:
 \[\begin{array}{c|c|c} n & 3^n\mod{169} & 3^n\mod{121}\\ \hline 0 & 1 & 1\\ 1 & 3 & 3\\ 2 & 9 & 9\\ 3 & 27 & 27\\ 4 & 81 & 81\\ 5 & 74 & 1\\ 6 & 53\\ 7 & 159\\ 8 & 139\\ 9 & 79\\ 10 & 68\\ 11 & 35\\ 12 & 105\\ 13 & 146\\ 14 & 100\\ 15 & 131\\ 16 & 55\\ 17 & 165\\ 18 & 157\\ 19 & 133\\ 20 & 61\\ 21 & 14\\ 22 & 42\\ 23 & 126\\ 24 & 40\\ 25 & 120\\ 26 & 22\\ 27 & 66\\ 28 & 29\\ 29 & 87\\ 30 & 92\\ 31 & 107\\ 32 & 152\\ 33 & 118\\ 34 & 16\\ 35 & 48\\ 36 & 144\\ 37 & 94\\ 38 & 113\\ 39 & 1\\ \end{array}\]

 The powers of $3$ repeat in cycles of $5$ and $39$ in modulo $121$ and modulo $169$, respectively. The answer is $\text{lcm}(5, 39) = \boxed{195}$.

 



\textbf{Solution 4 (Order+Bash)}

We have that 
 \[3^n \equiv 1 \pmod{143^2}.\]Now, $3^{110} \equiv 1 \pmod{11^2}$ so by the Fundamental Theorem of Orders, $\text{ord}_{11^2}(3)|110$ and with some bashing, we get that it is $5$. Similarly, we get that $\text{ord}_{13^2}(3)=39$. Now, $\text{lcm}(39,5)=\boxed{195}$ which is our desired solution.

 



\textbf{Solution 5 (Easy Binomial Theorem)}

We wish to find the smallest $n$ such that $3^n\equiv 1\pmod{143^2}$, so we want $n\equiv 1\pmod{121}$ and $n\equiv 1\pmod{169}$. Note that $3^5\equiv 1\pmod{121}$,  so $3^n$ repeats $121$ with a period of $5$, so $5|n$. Now, in order for $n\equiv 1\pmod{169}$, then $n\equiv 1\pmod{13}$. Because $3^3\equiv 1\pmod{13}$, $3^n$ repeats with a period of $3$, so $3|n$. Hence, we have that for some positive integer $p$, $3^n\equiv (3^3)^p\equiv (26+1)^p\equiv \binom{p}{0}26^p+\binom{p}{1}26^{p-1}....+\binom{p}{p-2}26^2+\binom{p}{p-1}26+\binom{p}{p}\equiv 26p+1\equiv 1\pmod{169}$, so $26p\equiv 0\pmod{169}$ and $p\equiv 0\pmod{13}$. Thus, we have that $5|n$, $3|n$, and $13|n$, so the smallest possible value of $n$ is $3\times5\times13=\boxed{195}$.-Stormersyle

 



\textbf{Solution 6 (LTE)}

We can see that $3^n-1 = 143^2*x$, which means that $v_{11}(3^n-1) \geq 2$, $v_{13}(3^n-1) \geq 2$. $v_{11}(3^n-1) = v_{11}(242) + v_{11}(\frac{n}{5})$, $v_{13}(3^n-1) = v_{13}(26) + v_{13}(\frac{n}{3})$ by the Lifting the Exponent lemma. From the first equation we gather that 5 divides n, while from the second equation we gather that both 13 and 3 divide n as $v_{13}(3^n-1) \geq 2$. Therefore the minimum possible value of n is $3\times5\times13=\boxed{195}$.

 -bradleyguo

 



\textbf{Solution 7 (LTE, more in-depth)}

As with the above solution, we are given $3^n \equiv 1 \pmod {143^2}.$ From CRT, we can separate these into. $3^n \equiv 1 \pmod {11^2}$ and $3^n \equiv 1 \pmod {13^2}.$

 (1) $3^n \equiv 1 \pmod {11^2}$Notice that for all $n \equiv 0 \pmod 5,$ the given equivalency holds. Hence, $5 \mid n.$

 (2) $3^n \equiv 1 \pmod {13^2}.$It is not exactly clear what $n$ would be here. We turn to LTE.

 First, note that $3^3 \equiv 1 \pmod {13} \Rightarrow \nu_{13}(3^3-1)=1.$

 For LTE, if we use $3^n-1^n,$ then we would need an odd prime to divide $3-1=2,$ absurd. If we instead use $3^{3n'}-1^{3n'}$ where $n=3n' \iff n \equiv 0 \pmod 3,$ then we can apply LTE.

 We have $\nu_{13}((3^3)^{n'}-(1^3)^{n'})=\nu_{13}(n')+\nu_{13}(3^3-1^3)=\nu_{13}(n')+1.$

 We need $\nu_{13}((3^3)^{n'}-(1^3)^{n'}) \geq 2 \Rightarrow \nu_{13}(n') \geq 1.$ Hence, $13 \mid n.$

 All in all, we need $3, 5, 13 \mid n,$ and the minimum such value is $195.$




\textbf{Solution 8}

Note that the problem is basically asking for the least positive integer $n$ such that $11^2 \cdot 13^2 | 3^n - 1.$ It is easy to see that $n = \text{lcm } (a, b),$ where $a$ is the least positive integer satisfying $11^2 | 3^a - 1$ and $b$ the least positive integer satisfying $13^2 | 3^b - 1$. Luckily, finding $a$ is a relatively trivial task, as one can simply notice that $3^5 = 243 \equiv 1 \mod 121$. However, finding $b$ is slightly more nontrivial. The order of $3^k$ modulo $13$ (which is $3$) is trivial to find, as one can either bash out a pattern of remainders upon dividing powers of $3$ by $13$, or one can notice that $3^3 = 27 \equiv 1 \mod 13$ (the latter which is the definition of period/orders by FLT). We can thus rewrite $3^3$ as $(2 \cdot 13 + 1) \mod 13^2$. Now suppose that 
 \[3^{3k} \equiv (13n + 1) \mod 13^2.\] I claim that $3^{3(k+1)} \equiv (13(n+2) + 1) \mod 13^2.$

 Proof:To find $3^{3(k+1)},$ we can simply multiply $3^{3k}$ by $3^3,$ which is congruent to $2 \cdot 13 + 1$ modulo $13^2$.By expanding the product out, we obtain 
 \[3^{3(k+1)} \equiv (13n + 1)(2 \cdot 13 + 1) = 13^2 \cdot 2n + 13n + 2 \cdot 13 + 1 \mod 13^2,\] and since the $13^2$ on the LHS cancels out, we're left with 
 \[13n + 2 \cdot 13 + 1 \mod 13^2 \implies 13(n+2) + 1 \mod 13^2\]. Thus, our claim is proven.Let $f(n)$ be the second to last digit when $3^{3k}$ is written in base $13^2$. Using our proof, it is easy to see that $f(n)$ satisfies the recurrence $f(1) = 2$ and $f(n+1) = f(n) + 2$. Since this implies $f(n) = 2n,$ we just have to find the least positive integer $n$ such that $2n$ is a multiple of $13$, which is trivially obtained as $13$.The least integer $n$ such that $3^n - 1$ is divisible by $13^2$ is $3 \cdot 13 = 39,$ so our final answer is $\text{lcm } (5, 39) = \boxed{195}.$

 -fidgetboss\_4000-minor edits made by srisainandan6

 



\textbf{Solution 9 (Official MAA)}

The requested positive integer is the least value of $n>0$ such that $3^n\equiv 1\pmod{143^2}.$ Note that $143=11\cdot 13.$ The least power of $3$ that is congruent to $1$ modulo $11^2$ is $3^5=243=2\cdot 11^2+1.$ It follows that. $3^n\equiv 1\pmod {11^2}$ if and only if $n=5j$ for some positive integer $j$.

 The least power of $3$ that is congruent to $1$ modulo $13$ is $3^3=27=2\cdot 13+1.$ It follows that $3^n\equiv 1\pmod{13}$ if and only if $n=3k$ for some positive integer $k$. Additionally, for some positive integer $k$, the Binomial Theorem shows that $3^{3k}=(26+1)^k=26\cdot k+1 \pmod{13^2}$. In particular, $3^n=3^{3k}\equiv 1\pmod {13^2}$ if and only if $k=13m$ for some positive integer $m$, that is, if and only if $n=39m.$

 Because $11^2$ and $13^2$ are relatively prime, $3^n\equiv 1\pmod {143^2}$ if and only if $3^n\equiv 1\pmod{11^2}$ and $3^n \equiv 1\pmod {13^2}$. This occurs if and only if $n$ is a multiple of both of the relatively prime integers $5$ and $39$, so the least possible value of $n$ is $5\cdot 39=\boxed{195}.$




\textbf{Solution 10 (Motivation and LTE)}

We first note that we wish to find $3^n \equiv 1 \pmod{11^2}$ and $3^n \equiv 1 \pmod{13^2}.$ Not thinking of anything else, we try a few numbers for the first condition to get that $5 \mid n.$ For the second condition, upon testing up to 729, we find that it doesn't yield a solution readily, so we use Lifting the Exponent from our toolkit to get that 
 \[v_{13}(27^m-1^m)=v_{13}(26)+v(m)=1+v(m), 3m = n\] which clearly implies $m=13$ and $39 | n.$ Our answer is then obviously $39 \cdot 5 = \boxed{195}.$




\textbf{Solution 11(LTE and a little bash)}

We write $3^n = 143^{n + 1} \cdot a_n + \cdots + 143^2 \cdot a_1 + 143^1 \cdot 0 + 143^0 \cdot 1 = 143^2(a_1 + \text{something}) + 1$ which immediately tells us that $3^n \equiv 1 \mod (143^2)$. Notice $143^2 = 11^2 \cdot 13^2$. Thus, we solve $3^n \equiv 1 \mod 11$ first. We see that $3^2 = 9 \equiv -2 \mod 11$ so squaring gives $3^4 \equiv 4 \mod 11$ and we instantly see that multiplying by $3$ gives $3^5 \equiv 12 \mod 11 \equiv 1 \mod 11$ thus $5 | n$. Now, we let $n = 5k$. We actually want $3^n \equiv 1 \mod (11^2)$. Then we get

 $3^{5k} = (3^5)^k = 243^k = (2 \cdot 11^2 + 1)^k \equiv 1^k \equiv 1 \mod 11^2$ so $n = 5k$ automatically satisfies the $11^2$ condition. Now, for $3^n \equiv 1 \mod 13$, we instantly see $3^3 = 27 \equiv 1 \mod 13$ so $3 | n$. Thus, we can write $n = 5 \cdot 3 \cdot k = 15k$. Thus, we need

 $3^{15k} \equiv 1 \mod 169$. Now computing $3^{15}$ requires some bash, but it isn't bad at all after this is computed. Ok, first $3^15 = (3^5)^3 = (243)^3 \equiv (243 - 169)^3 \mod 169 \equiv 74^3 \mod 169$. We have to manually compute $74^2 = (70 + 4)^2 = 4900 + 560 + 16 = 5476$ and we do long division to get $5476 \equiv 68 \mod 169$ thus we compute $68 \cdot 74 \mod 169$ and notice $68 \cdot 74 = (71 - 3)(71 + 3) = (70 + 1)^2 - 9 = 4900 + 140 + 1 - 9 = 5040 - 8 = 5032$ thus $5032 \mod 169 \equiv 131 \mod 169$ from long division. Thus, we just need to solve

 $131^k \equiv 1 \mod 169$. Here, we use Lifting the Exponent Lemma or LTE. We have $v_{13}(131^k - 1^k) = 2$ since $169$ has $2$ powers of $13$ and $v_x(y)$ denotes the highest power of $x$ that divides $y$. By LTE, since $13 | 131 - 1$, we have

 $v_{13}(131 - 1) + v_{13}(k) = v_{13}(131^k - 1^k) = 2 \implies v_{13}(13 \cdot 10) + v_{13}(k) = 2 \implies 1 + v_{13}(k) = 2 \implies v_{13}(k) = 1$ thus $13 | k$.

 All in all, we see that we can now write $k = 13j$ so $n = 15k = 15 \cdot 13j = 195j$ and this value satisfies all the conditions of the problem. The minimum positive integer value of $n$ occurs at $j = 1$ thus $\boxed{195}$



	\item (\textit{AIME I 2018}) For every subset $T$ of $U = \{ 1,2,3,\ldots,18 \}$, let $s(T)$ be the sum of the elements of $T$, with $s(\emptyset)$ defined to be $0$. If $T$ is chosen at random among all subsets of $U$, the probability that $s(T)$ is divisible by $3$ is $\frac{m}{n}$, where $m$ and $n$ are relatively prime positive integers. Find $m$.


 
 \textbf{Answer: }683



\textbf{Solution 1}

The question asks us for the probability that a randomly chosen subset of the set of the first 18 positive integers has the property that the sum of its elements is divisible by 3. Note that the total number of subsets is $2^{18}$ because each element can either be in or not in the subset. To find the probability, we will find the total numbers of ways the problem can occur and divide by $2^{18}$.

 To simplify the problem, let’s convert the set to mod 3:

 
 \[U' = \{1,2,0,1,2,0,1,2,0,1,2,0,1,2,0,1,2,0\}\]

 Note that there are six elements congruent to 0 mod 3, 6 congruent to 1 mod 3, and 6 congruent to 2 mod 3. After conversion to mod three, the problem is the same but we’re dealing with much simpler numbers. Let’s apply casework on this converted set based on $S = s(T')$, the sum of the elements of a subset $T'$ of $U'$.

 $\textbf{Case 1: }S=0$

 In this case, we can restrict the subsets to subsets that only contain 0. There are six 0’s and each one can be in or out of the subset, for a total of $2^{6} = 64$ subsets. In fact, for every case we will have, we can always add a sequence of 0’s and the total sum will not change, so we will have to multiply by 64 for each case. Therefore, let’s just calculate the total number of ways we can have each case and remember to multiply it in after summing up the cases. This is equivalent to finding the number of ways you can choose such subsets without including the 0's. Therefore this case gives us $\boxed{1}$ way.

 $\textbf{Case 2: }S= 3$

 In this case and each of the subsequent cases, we denote the number of 1’s in the case and the number of two’s in the case as $x, y$ respectively. Then in this case we have two subcases;

 $x, y = 3,0:$ This case has $\tbinom{6}{3} \cdot \tbinom{6}{0} = 20$ ways.

 $x, y = 1,1:$ This case has $\tbinom{6}{1} \cdot \tbinom{6}{1} = 36$ ways.

 In total, this case has $20+36=\boxed{56}$ ways.

 $\textbf{Case 3: }S=6$

 In this case, there are 4 subcases;

 $x, y = 6,0:$ This case has $\tbinom{6}{6} \cdot \tbinom{6}{0} = 1$ way.

 $x, y = 4,1:$ This case has $\tbinom{6}{4} \cdot \tbinom{6}{1} = 90$ ways.

 $x, y = 2,2:$ This case has $\tbinom{6}{2} \cdot \tbinom{6}{2} = 225$ ways.

 $x, y = 0,3:$ This case has $\tbinom{6}{0} \cdot \tbinom{6}{3} = 20$ ways.

 In total, this case has $1+90+225+20=\boxed{336}$ ways.

 Note that for each case, the \# of 1’s goes down by 2 and the \# of 2’s goes up by 1. This is because the sum is fixed, so when we change one it must be balanced out by the other. This is similar to the Diophantine equation $x + 2y= S$ and the total number of solutions will be $\tbinom{6}{x} \cdot \tbinom{6}{y}$. From here we continue our casework, and our subcases will be coming out quickly as we have realized this relation.

 $\textbf{Case 4: }S=9$

 In this case we have 3 subcases:

 $x, y = 5,2:$ This case has $\tbinom{6}{5} \cdot \tbinom{6}{1} = 90$ ways.

 $x, y = 3,3:$ This case has $\tbinom{6}{3} \cdot \tbinom{6}{3} = 400$ ways.

 $x, y = 1,4:$ This case has $\tbinom{6}{1} \cdot \tbinom{6}{4} = 90$ ways.

 In total, this case has $90+400+90=\boxed{580}$ ways.

 $\textbf{Case 5: } S=12$

 In this case we have 4 subcases:

 $x, y = 6,3:$ This case has $\tbinom{6}{6} \cdot \tbinom{6}{3} = 20$ ways.

 $x, y = 4,4:$ This case has $\tbinom{6}{4} \cdot \tbinom{6}{4} = 225$ ways.

 $x, y = 2,5:$ This case has $\tbinom{6}{2} \cdot \tbinom{6}{5} = 90$ ways.

 $x, y = 0,6:$ This case has $\tbinom{6}{0} \cdot \tbinom{6}{6} = 1$ way.

 Therefore the total number of ways in this case is $20 + 225 + 90 + 1=\boxed{336}$. Here we notice something interesting. Not only is the answer the same as Case 3, the subcases deliver the exact same answers, just in reverse order. Why is that?

 We discover the pattern that the values of $x, y$ of each subcase of Case 5 can be obtained by subtracting the corresponding values of $x, y$ of each subcase in Case 3 from 6 ( For example, the subcase 0, 6 in Case 5 corresponds to the subcase 6, 0 in Case 3). Then by the combinatorial identity, $\tbinom{6}{k} = \tbinom{6}{6-k}$, which is why each subcase evaluates to the same number. But can we extend this reasoning to all subcases to save time?

 Let’s write this proof formally. Let $W_S$ be the number of subsets of the set $\{1,2,1,2,1,2,1,2,1,2,1,2\}$ (where each 1 and 2 is distinguishable) such that the sum of the elements of the subset is $S$ and $S$ is divisible by 3. We define the empty set as having a sum of 0. We claim that $W_S = W_{18-S}$.

 $W_S = \sum_{i=1}^{|D|} \tbinom{6}{x_i}\tbinom{6}{y_i}$ if and only if there exists $x_i, y_i$ that satisfy $0\leq x_i \leq 6$, $0\leq y_i \leq 6$, $x_i + 2y_i = S$, and $D$  is the set of the pairs $x_i, y_i$. This is because for each pair $x_i$, $y_i$ there are six elements of the same residue mod(3) to choose $x_i$ and $y_i$ numbers from, and their value sum must be $S$.

 Let all $x_n$, $y_n$ satisfy $x_n = 6-x_i$ and $y_n = 6-y_i$.We can rewrite the equation $x_i+ 2y_i = S \implies -x_i- 2y_i = -S \implies 6-x_i+ 6-2y_i= 12 - S$$\implies 6-x_i+12-2y_i = 18-S \implies 6-x_i + 2(6-y_i) = 18-S$
 \[\implies x_n + 2y_n = 18 - S\]

 Therefore, since $0\leq x_i, y_i\leq 6$ and $x_n = 6-x_i$ and $y_n = 6-y_i$, $0\leq x_n, y_n\leq 6$. As shown above, $x_n + 2y_n = 18 - S$ and since $S$ and 18 are divisible by 3, 18 -$S$ is divisible by 3. Therefore, by the fact that $W_S = \sum_{i=1}^{|D|} \tbinom{6}{x_i}\tbinom{6}{y_i}$, we have that;

 $W_{18-S} = \sum_{n=1}^{|D|} \tbinom{6}{x_n}\tbinom{6}{y_n} \implies W_{18-S} =  \sum_{i=1}^{|D|} \tbinom{6}{6-x_i}\tbinom{6}{6-y_i} \implies W_{18-S} =  \sum_{i=1}^{|D|} \tbinom{6}{x_i}\tbinom{6}{y_i} = W_S$, proving our claim.

 We have found our shortcut, so instead of bashing out the remaining cases, we can use this symmetry. The total number of ways over all the cases is $\sum_{k=0}^{6} W_{3k} = 2 \cdot (1+56+336)+580 = 1366$. The final answer is $\frac{2^{6}\cdot 1366}{2^{18}} = \frac{1366}{2^{12}} = \frac{683}{2^{11}}.$

 There are no more 2’s left to factor out of the numerator, so we are done and the answer is $\boxed{\boxed{683}}$.

 



\textbf{Solution 2}

Consider the elements of $U$ modulo $3.$

 Ignore the $0$'s because we're gonna multiply $\binom{6}{0}+..+\binom{6}{6}=2^6$ at the end. Let $a$ be the $1's$ and $b$ be the $2's$. The key here is that $2 \equiv -1 \pmod{3}$ so the difference between the number of $a$ and $b$ is a multiple of $3$.

 1. Counted twice because $a$ and $b$ can be switched:

 $6a$
$6a,3b$
$5a,2b$
$4a,b$
$3a$

 2. Counted once:

 $6a,6b$...$0a,0b$

 By \href{https://artofproblemsolving.com/wiki/index.php?title=Combinatorial\_identity}{ Vandermonde's Identity} on the second case, this is $2^6 \left( 2\left(1+20+90+90+20\right) + \binom{12}{6} \right)\implies \boxed{683}$




\textbf{Solution 3 (Triple Vandermonde's)}

Elements $0 \pmod{3}$ can either be included or excluded, for a total of $2^6$. Then, like the previous solution, let $a$ be the number of elements $1 \pmod{3}$ and $b$ be the number of elements $2 \pmod{3}$. Then, $a + 2b \equiv 0 \pmod{3} \implies a - b \equiv 0 \pmod{3}$. Since $0 \le a, b \le 6$, there are 3 cases:

 $\>\>\>\> 1. \> a - b = 0$

 Then, we have,
 \[\binom{6}{0}\binom{6}{0} + \binom{6}{1}\binom{6}{1} + \cdots + \binom{6}{6} \binom{6}{6}\]

 Using the substitution $\binom{6}{k} = \binom{6}{6-k}$ on every second binomial coefficient, it is clear that the bottom numbers sum to $6$, and the ones above sum to $12$. Apply Vandermonde's, we obtain $\binom{12}{6}$
$\>\>\>\> 2. \> a - b = \pm 3$

 For $a - b = 3$, using the same substitution and applying Vandermonde's, we get:
 \[\binom{6}{3} \binom{6}{0} + \binom{6}{4} \binom{6}{1} + \cdots + \binom{6}{6} \binom{6}{3} = \binom{12}{3}\]

 Then, for $a - b = -3$, it is completely symmetric with $a - b = 3$. We get $\binom{12}{3}$
$\>\>\>\> 3. \> a - b = \pm 6$

 We have $\binom{6}{6} \binom{6}{0} + \binom{6}{0} \binom{6}{6} = 2$

 We get $\frac{2^6 (\binom{12}{6} + 2 \binom{12}{3} + 2)}{2^{18}} = \frac{683}{2^{11}} \implies m = \boxed{683}$.

 



\textbf{Solution 4 (quickest recursion)}

We generalize the problem to the set $U_n = \{1, 2, 3, \cdots{}, 3n\}.$ 

 Let $T_n$ be a subset taken from $U_n$, and denote by $A_n, B_n,$ and $C_n$ the probabilities that $s(T_n)$ is congruent to $0, 1,$ and $2 \pmod{3}$ respectively. Computing for $n=1$, we have $A_1 = \frac{1}{2}, B_1 = C_1 = \frac{1}{4}.$ Tracking ways to make each possible residue based on the previous residue of $s(T_n),$ we obtain the recurrence

 
 \[A_{n+1} = \frac{B_n}{4} + \frac{C_n}{4} + \frac{A_n}{2} = \frac{2A+B+C}{4} = \frac{1+A}{4}\]where the last equality follows from the fact that $A_n+B_n+C_n = 1.$

 Solving this recurrence (or simply manually computing upwards) yields $A_6 = \frac{683}{2048} \implies \boxed{683}$.

 



\textbf{Solution 5 (Symmetry)}

Note that in general, the answer will be around $1/3$ of $2^{18}$. We will use this to our advantage. Partition $U$ into disjoint subsets $\{1,2,3\}, \{4,5,6\}, \cdots, \{16,17,18\}$. Within each of these subsets, there are 3 possible remainders$\mod 3$ depending on what elements we choose to include into $T$: (Using set $\{1,2,3\}$ for demonstration purposes) 

 Remainder of zero: Include $\{3\}, \{1,2\}, \{\}, \{1,2,3\}$

 Remainder of one: Include $\{1\}, \{1,3\}$

 Remainder of two: Include $\{2\}, \{2,3\}$

 Suppose all subsets are of the form $\{\}, \{1,2,3\}$. Then, since both choices are $0 \pmod{3}$, we can choose either one for all six subsets, giving us $2^6$ combinations.

 On the other hand, suppose there exists a subset $V$ that isn't of the form $\{\}, \{1,2,3\}$. Then, for $V$, it is equally likely that a remainder of zero, one, or two is chosen, since each remainder has $2$ ways to be achieved, i.e. there is a $1/3$ chance for each. Consider the sum of the other subsets $T \setminus V$. The sum is either $0, 1, 2 \pmod{3}$. Whatever that remainder $r$ might be, we can always choose $3-r$ as the remainder for our set $V$, giving us a total of $r + 3 - r \equiv 0 \pmod{3}$. The probability we choose remainder $3-r$ is $1/3$ as previously mentioned. So, we get $\frac{1}{3}(2^{18} - 2^6)$ total combinations.

 Therefore, we get $\frac{2^6 + \frac{1}{3}(2^{18} - 2^6)}{2^{18}} = \frac{683}{2^{11}} \implies m = \boxed{683}$




\textbf{Solution 6}

Rewrite the set after mod 3 as above

 1 2 0 1 2 0 1 2 0 1 2 0 1 2 0 1 2 0

 All 0s can be omitted 

 Case 1No 1 No 2$1$

 Case 2$222$$20$

 Case 3$222222$$1$

 Case 4$12$$6*6=36$

 Case 5$12222$$6*15=90$

 Case 6$1122$$15*15=225$

 Case 7$1122222$$15*6=90$

 Case 8$111$$20$

 Case 9$111222$$20*20=400$

 Case 10$111222222$$20$

 Case 11$11112$$15*6=90$

 Case 12$11112222$$15*15=225$

 Case 13$1111122$$6*15=90$

 Case 14$1111122222$$6*6=36$

 Case 15$111111$$1$

 Case 16$111111222$$20$

 Case 17$111111222222$$1$

 Total $1+20+1+36+90+225+90+20+400+20+90+225+90+36+1+20+1=484+360+450+72=1366$
$P=\frac{1366}{2^{12}}=\frac{683}{2^{11}}$

 ANS=$\boxed{683}$

 By SpecialBeing2017

 



\textbf{Solution 7}

Consider the numbers $\{1,4,7,10,13,16\}$. Each of those are congruent to $1 \pmod 3$. There is ${6 \choose 0}=1$ way to choose zero numbers ${6 \choose 1}=6$ ways to choose $1$ and so on. There ends up being  ${6 \choose 0}+{6 \choose 3}+{6 \choose 6} = 22$ possible subsets congruent to $0\pmod 3$.  There are $2^6=64$ possible subsets of these numbers. By symmetry there are $21$ subsets each for $1 \pmod 3$ and $2 \pmod3$.

 We get the same numbers for the subsets of $\{2,5,8,11,14,17\}$.

 For $\{3,6,9,12,15,18\}$, all $2^6$ subsets are $0 \pmod3$.

 So the probability is: $\frac{(22\cdot22+21\cdot21+21\cdot21)\cdot2^6}{2^{18}}=\frac{683}{2^{11}}$




\textbf{Solution 8}

Notice that six numbers are $0\pmod3$, six are $1\pmod3$, and six are $2\pmod3$.  Having numbers $0\pmod3$ will not change the remainder when $s(T)$ is divided by $3$, so we can choose any number of these in our subset.  We ignore these for now.  The number of numbers that are $1\pmod3$, minus the number of numbers that are $2\pmod3$, must be a multiple of $3$, possibly zero or negative.  We can now split into cases based on how many numbers that are $1\pmod3$ are in the set.

 CASE 1- $0$, $3$, or $6$ integers:  There can be $0$, $3$, or $6$ integers that are $2\pmod3$.  We can choose these in 
 \[\left(\binom60+\binom63+\binom66\right)\cdot\left(\binom60+\binom63+\binom66\right)=(1+20+1)^2=484.\]

 CASE 2- $1$ or $4$ integers:  There can be $2$ or $5$ integers that are $2\pmod3$.  We can choose these in 
 \[\left(\binom61+\binom64\right)\cdot\left(\binom62+\binom65\right)=(6+15)^2=441.\]

 CASE 3- $2$ or $5$ integers:  There can be $1$ or $4$ integers that are $2\pmod3$.  We can choose these in 
 \[\left(\binom62+\binom65\right)\cdot\left(\binom61+\binom64\right)=(15+6)^2=441.\]

 Adding these up, we get that there are $1366$ ways to choose the numbers such that their sum is a multiple of three.  Putting back in the possibility that there can be multiples of $3$ in our set, we have that there are 
 \[1366\cdot\left(\binom60+\binom61+\binom62+\binom63+\binom64+\binom65+\binom66\right)=1366\cdot2^6\] subsets $T$ with a sum that is a multiple of $3$.  Since there are $2^{18}$ total subsets, the probability is 
 \[\frac{1366\cdot2^6}{2^{18}}=\frac{683}{2^{11}},\] so the answer is $\boxed{683}$.

 



\textbf{Solution 9 (Generating Functions)}

We use generating functions. Each element of $U$ has two choices that occur with equal probability--either it is in the set or out of the set. Therefore, given $n\in U$, the probability generating function is
 \[\frac{1}{2}+\frac{1}{2}x^n.\]Therefore, in the generating function
 \[\frac{1}{2^{18}}(1+x)(1+x^2)(1+x^3)\cdots (1+x^{18}),\]the coefficient of $x^k$ represents the probability of obtaining a sum of $k$. We wish to find the sum of the coefficients of all terms of the form $x^{3k}$ (where $k \in \mathbb{Z}_{\geq0}$). If $\omega=e^{2\pi i/3}$ is a cube root of unity, then it is well know that for a polynomial $P(x)$, 
 \[\frac{P(1)+P(\omega)+P(\omega^2)}{3}\]will yield the sum of the coefficients of the terms of the form $x^{3k}$. (This is known as the third Roots of Unity Filter.) Then we find
 \begin{align*} \frac{1}{2^{18}}(1+1)(1+1^2)(1+1^3)\cdots (1+1^{18})&=1\\\frac{1}{2^{18}}(1+\omega)(1+\omega^2)(1+\omega^3)\cdots (1+\omega^{18})&=\frac{1}{2^{12}}\\\frac{1}{2^{18}}(1+\omega^2)(1+\omega^4)(1+\omega^6)\cdots (1+\omega^{36})&=\frac{1}{2^{12}}. \end{align*}
To evaluate the last two products, we utilized the facts that $\omega^3=1$ and $1+\omega+\omega^2=0$. Therefore, the desired probability is
 \[\frac{1+1/2^{12}+1/2^{12}}{3}=\frac{683}{2^{11}}.\]Thus the answer is $\boxed{683}$.

 



\textbf{Solution 10}

Define a set that "works" to be a set for which the sum of the terms is $0$ mod $3$. The given set mod $3$ is
 \[\{1,2,0,1,2,0,1,2,0,1,2,0,1,2,0,1,2,0\}.\]Let $w(N)$ be the number of subsets of the first $N$ terms of this set that "work."Consider just the first $3$ terms:
 \[\{1,2,0\}.\]There are $2^3=8$ total subsets, and $w(3)=4$ (the subsets are $\emptyset, \{0\}, \{1,2\},$ and $\{1,2,0\}$). Now consider the first $6$ terms:
 \[\{1,2,0,1,2,0\}.\]To find $w(6)$, we set aside the last $3$ terms as a "reserve" that we can pair with subsets of the first $3$ terms (which we looked at earlier). 

 First, all $2^3$ subsets of the first $3$ terms can either be paired with a $1$ or a $2$ (or nothing) from the "reserve" terms so that they "work," creating $2^3$ subsets that "work" already. But for each of these, we have the option to add a $0$ from the reserve, so we now have $2\cdot2^3$ subsets that "work." For each of the $w(3)=4$ subsets of the first $3$ terms that "work", we can also add on a $\{1,2\}$ or a $\{1,2,0\}$ from the reserves, creating $2w(3)$ new subsets that "work." And that covers all of them. With all of this information, we can write $w(6)$ as
 \[w(6)=2\cdot2^3+2w(3)=2(2^3+w(3)).\]The same process can be used to calculate $w(9)$ then $w(12)$ etc. so we can generalize it to
 \[w(N)=2(2^{N-3}+w(N-3)).\]Thus, we calculate $w(18)$ with this formula:
 \[w(18)=2( 2^{15} + 2( 2^{12} + 2( 2^9 + 2( 2^6 + 2( 2^3 + 4 ) ) ) ) )=87424.\]Because there are $2^{18}$ total possible subsets, the desired probability is
 \[\frac{w(18)}{2^{18}}=\frac{87424}{2^{18}}=\frac{683}{2048},\]so the answer is $\boxed{683}.$




\textbf{Solution 11}

Try smaller cases and find a pattern. Using similar casework as in Solution 1, we can easily find the desired probability if $U$ is of a small size. 

 If $U = \{ 1,2,0\} \pmod 3$, then $4$ out of $8$ subsets work, for a probability of $\frac12$.

 If $U = \{ 1,2,0,1,2,0\} \pmod 3$, then $24$ out of $64$ subsets work, for a probability of $\frac38$.

 If $U = \{ 1,2,0,1,2,0,1,2,0\} \pmod 3$, then $176$ out of $512$ subsets work, for a probability of $\frac{11}{32}$.

 Let $a_n$ be the numerator of the desired probability if $U$ is of size $3n$. Then $a_1 = 1, a_2 = 3,$ and $a_3 = 11$. Noticing that the denominators are multiplied by $4$ each time, we can conclude that the pattern of the numerators is $a_n = 4a_{n-1} - 1$, so $a_6 = \boxed{683}$.

 



\textbf{Solution 12 (Quick guesswork for about 2 minutes remaining)}

We conjecture that the difference from the probability will be as small as possible from $\frac{1}{3}$ (The value approached as $n \rightarrow \infty$ , where $n$ is the number of terms in the largest subset).

 
We also see that there are $2^{18}$ subsets and know the denominator will be a power of $2$ (since the numerator is an integer).

 We essentially want to guess that the greatest integer $n$ satisfying $\left(\frac{2^n}{3}\right)-1 <1000$ can be plugged in to get the solution of round $\left( \frac{2^n}{3} \right)$.

 
We see that this occurs at $n=11$, and get round $\left( \frac{2^{11}}{3} \right)=$round($682.66...$)= $683$.

 



\textbf{Solution 13 (Very quick recursion)}

The total number of subsets is simply $2^{18}.$ Now we need to find the number of subsets that have a sum divisible by $3.$

 Ignore the 6 numbers in the list that are divisible by 3. We look only at the number of subsets of $\{1, 2, 4, 5, 7, 8, 10, 11, 13, 14, 16, 17\}$ then multiply by $2^6$ at the end. This is because adding a multiple of $3$ to the sum will not change whether it is divisible by $3$ or not and for each of the $6$ multiples, we have two options of whether to add it to the subset or not. 

 Now, let $f(3k)$ be the number of subsets of $\{1, 2, 4, 5, \dots, 3k-2, 3k-1\}$ that have a sum divisible by $3$. There are $2k$ numbers in the set. Let us look at the first $2k - 2$ numbers: all subsets of the first $2k - 2$ numbers will have a residue $0, 1,$ or $2$ mod $3.$ If it is $0,$ add both $3k-2$ and $3k-1$ to the subset. If it is $1$ add just $3k-1$ to the set, and if it is $2$ add just $3k-2.$  We don't have to compute all three cases separately; since there is a 1 to 1 correspondence, we only need the sum of the three cases (which we know is $2^{2k-2}$). 

 Now, this counts all the subsets except for those that include neither $3k-1$ nor $3k-2.$ However, this is just $f(3k - 3).$ Thus, $f(3k) = 2^{2k-2} + f(3k-3)$ and the base case is $f(0) = 1.$ We have, 

 $f(18) = 2^{10} + f(15) = 2^{10} + 2^{8} + f(12) = \dots = 2^{10} + 2^8 + 2^6 + 2^4 + 2^2 + 2^0 + f(0).$$=2^{10} + 2^8 + \dots + 2^2 + 2.$

 Multiplying this by $2^6$ we have, 
 \[\frac{2^7(1 + 2^1 + 2^3 + 2^5 + 2^7 + 2^9)}{2^{18}} .\]The $2^7$ cancels out with the denominator. However, the second term in the product is obviously odd and so does not cancel further with the denominator, which is just a power of $2.$ Thus, we want to find $1 + 2 + 2^3 + 2^5 + 2^7 + 2^9,$ which equals $\boxed{683}.$ 

 Note: While the explanation may seem a bit complicated, the concept behind it is pretty simple. And once you get the solution, computing the answer doesn't take much time.

 


 



\textbf{Solution 14}

The total number of subsets is $\sum_{i=0}^{18}\tbinom{18}{i}=2^{18}$. If $s(T)\equiv 0\bmod{3}$, the sum of the elements divides 3. We can rewrite the set as 6 0s, 6 1s, and 6 2s. We can ignore the zeros for now, since they won't influence the sum so we focus on each configuration of the 6 1s and 6 2s such that the sum is divisible by 3 and then multiply by $\sum_{i=0}^{6}\tbinom{6}{i}=2^6$ at the end. We proceed with casework on the number of values that are equivalent to $2\pmod{3}$ as follows:

 Case 1: There are zero elements:

 Then, there are only configurations of the values congruent to 1 mod 3. There are $\tbinom{6}{0}+\tbinom{6}{3}+\tbinom{6}{6}$ ways to assign values in the set that are congruent to 1 mod 3 such that the sum of those values divides 3. Therefore there are $22$ configurations for this case.

 Case 2: There is one element:

 There are $6$ ways to choose which element is included in the subset and $\tbinom{6}{1}+\tbinom{6}{4}$ ways to assign values to the numbers congruent to 1 mod 3 in the set. Therefore there are $6\cdot21=126$ configurations for this case.

 Case 3: There are two elements:

 There are $\tbinom{6}{2}=12$ ways to choose the 2 elements in the set that are congruent to 2 mod 3 and $\tbinom{6}{2}+\tbinom{6}{5}$ possible ways to assign values to the numbers congruent to 1 mod 3. Therefore there are $15\cdot21=315$ configurations for this case.

 Case 4: There are three elements:

 There are $\tbinom{6}{3}=20$ ways to choose three elements in the set that are congruent to 2 mod 3. There are $\tbinom{6}{0}+\tbinom{6}{3}+\tbinom{6}{6}$ ways to assign values to the numbers congruent to 1 mod 3. Therefore, there are $20\cdot22=440$ configurations for this case.

 Case 5, Case 6, and Case 7 are symmetric to Case 3, Case 2, and Case 1 respectively so we can multiply by 2 for those cases. 

 Putting all the cases together we obtain $2(22+126+315) + 440=1366$. Multiplying by $2^6$ gives $2^6\cdot1366=2^7\cdot683$. Since $n=2^{18}$, $\frac{m}{n}=\frac{683}{2^{11}}$. Therefore, $m=\boxed{683}$



	\item (\textit{AIME I 2018}) Let $\triangle ABC$ have side lengths $AB=30$, $BC=32$, and $AC=34$. Point $X$ lies in the interior of $\overline{BC}$, and points $I_1$ and $I_2$ are the incenters of $\triangle ABX$ and $\triangle ACX$, respectively. Find the minimum possible area of $\triangle AI_1I_2$ as $X$ varies along $\overline{BC}$.


 
 \textbf{Answer: }126



\textbf{Solution 1 (Official MAA)}

First note that 
 \[\angle I_1AI_2 = \angle I_1AX + \angle XAI_2 = \frac{\angle BAX}2 + \frac{\angle CAX}2 = \frac{\angle A}2\] is a constant not depending on $X$, so by $[AI_1I_2] = \frac12(AI_1)(AI_2)\sin\angle I_1AI_2$ it suffices to minimize $(AI_1)(AI_2)$.  Let $a = BC$, $b = AC$, $c = AB$, and $\alpha = \angle AXB$.  Remark that 
 \[\angle AI_1B = 180^\circ - (\angle I_1AB + \angle I_1BA) = 180^\circ - \frac12(180^\circ - \alpha) = 90^\circ + \frac\alpha 2.\] Applying the Law of Sines to $\triangle ABI_1$ gives 
 \[\frac{AI_1}{AB} = \frac{\sin\angle ABI_1}{\sin\angle AI_1B}\qquad\Rightarrow\qquad AI_1 = \frac{c\sin\frac B2}{\cos\frac\alpha 2}.\] Analogously one can derive $AI_2 = \frac{b\sin\frac C2}{\sin\frac\alpha 2}$, and so 
 \[[AI_1I_2] = \frac{bc\sin\frac A2 \sin\frac B2\sin\frac C2}{2\cos\frac\alpha 2\sin\frac\alpha 2} = \frac{bc\sin\frac A2 \sin\frac B2\sin\frac C2}{\sin\alpha}\geq bc\sin\frac A2 \sin\frac B2\sin\frac C2,\] with equality when $\alpha = 90^\circ$, that is, when $X$ is the foot of the perpendicular from $A$ to $\overline{BC}$.  In this case the desired area is $bc\sin\frac A2\sin\frac B2\sin\frac C2$.  To make this feasible to compute, note that 
 \[\sin\frac A2=\sqrt{\frac{1-\cos A}2}=\sqrt{\frac{1-\frac{b^2+c^2-a^2}{2bc}}2} = \sqrt{\dfrac{(a-b+c)(a+b-c)}{4bc}}.\] Applying similar logic to $\sin \frac B2$ and $\sin\frac C2$ and simplifying yields a final answer of 
 \begin{align*}bc\sin\frac A2\sin\frac B2\sin\frac C2&=bc\cdot\dfrac{(a-b+c)(b-c+a)(c-a+b)}{8abc}\\&=\dfrac{(30-32+34)(32-34+30)(34-30+32)}{8\cdot 32}=\boxed{126}.\end{align*}


            \begin{itemize}
        \item Notice that we truly did minimize the area for $[A I_1 I_2]$ because $b, c, \angle A, \angle B, \angle C$ are all constants while only $\sin \alpha$ is variable, so maximizing $\sin \alpha$ would minimize the area.
            \end{itemize}
            



\textbf{Solution 2 (Similar to Official MAA)}

It's clear that $\angle I_{1}AI_{2}=\frac{1}{2}\angle BAX+\frac{1}{2}\angle CAX=\frac{1}{2}\angle A$. Thus 
 \begin{align*} AI_{1}I_{2}&=\frac{1}{2}\cdot AI_{1}\cdot AI_{2}\cdot\sin\angle I_{1}AI_{2} \\ &=\frac{1}{2}\cdot AI_{1}\cdot AI_{2}\cdot\sin\left(\frac{1}{2}\angle A\right).\end{align*}
 By the Law of Sines on $\triangle AI_{1}B$, 
 \[\frac{AI_{1}}{\sin\left(\frac{1}{2}\angle B\right)}=\frac{AB}{\sin\angle AI_{1}B}.\] Similarly, 
 \[\frac{AI_{2}}{\sin\left(\frac{1}{2}\angle C\right)}=\frac{AC}{\sin\angle AI_{2}C}.\] It is well known that 
 \[\angle AI_{1}B=90+\frac{1}{2}\angle AXB~~~\text{and}~~~\angle AI_{2}C=90+\frac{1}{2}\angle AXC.\] Denote $\alpha=\frac{1}{2}\angle AXB$ and $\theta=\frac{1}{2}\angle AXC$, with $\alpha+\theta=90^{\circ}$. Thus $\sin\alpha=\cos\theta$ and 
 \begin{align*}\frac{AI_{1}}{\sin\left(\frac{1}{2}\angle B\right)}=\frac{AB}{\sin\angle AI_{1}B}\Longrightarrow\frac{AB}{\sin\left(90^{\circ}+\alpha\right)}\Longrightarrow\frac{AB}{\cos\alpha} \\\frac{AI_{2}}{\sin\left(\frac{1}{2}\angle C\right)}=\frac{AC}{\sin\angle AI_{2}C}\Longrightarrow\frac{AC}{\sin\left(90^{\circ}+\theta\right)}\Longrightarrow\frac{AC}{\cos\theta}\Longrightarrow\frac{AC}{\sin\alpha}.\end{align*}
 Thus 
 \[AI_{1}=\frac{AB\sin\left(\frac{1}{2}\angle B\right)}{\cos\alpha}~~~\text{and}~~~AI_{2}=\frac{AC\sin\left(\frac{1}{2}\angle C\right)}{\sin\alpha}\] so 
 \begin{align*}[AI_{1}I_{2}]&=\frac{1}{2}\cdot AI_{1}\cdot AI_{2}\cdot\sin\left(\frac{1}{2}\angle A\right) \\ &=\frac{AB\sin\left(\frac{1}{2}\angle B\right)\cdot AC\sin\left(\frac{1}{2}\angle C\right)\cdot\sin\left(\frac{1}{2}\angle A\right)}{2\sin\alpha\cos\alpha} \\ &=\frac{AB\sin\left(\frac{1}{2}\angle B\right)\cdot AC\sin\left(\frac{1}{2}\angle C\right)\cdot\sin\left(\frac{1}{2}\angle A\right)}{\sin(2\alpha)}.\end{align*}
 We intend to minimize this expression, which is equivalent to maximizing $\sin(2\alpha)$, and that occurs when $\alpha=45^{\circ}$, or $\angle AXB=90^{\circ}$. Ergo, $X$ is the foot of the altitude from $A$ to $\overline{BC}$. In that case, we intend to compute 
 \[AB\cdot AC\cdot\sin\left(\frac{1}{2}\angle B\right)\cdot\sin\left(\frac{1}{2}\angle C\right)\cdot\sin\left(\frac{1}{2}\angle A\right).\] Recall that 
 \[\sin\left(\frac{1}{2}\angle B\right)=\sqrt{\frac{1-\cos\angle B}{2}}\] and similarly for angles $C$ and $A$. Applying the Law of Cosines to each angle of $\triangle ABC$ gives 
 \begin{align*}\angle B&:\cos\angle B=\frac{30^{2}+32^{2}-34^{2}}{2\cdot 30\cdot 32}=\frac{2}{5} \\ \angle C&:\cos\angle C=\frac{32^{2}+34^{2}-30^{2}}{2\cdot 32\cdot 34}=\frac{10}{17} \\ \angle A&:\cos\angle A=\frac{30^{2}+34^{2}-32^{2}}{2\cdot 30\cdot 34}=\frac{43}{85}.\end{align*}
 Thus 
 \begin{align*}\sin\left(\frac{1}{2}\angle B\right)=\sqrt{\frac{1-\frac{2}{5}}{2}}=\sqrt{\frac{3}{10}} \\ \sin\left(\frac{1}{2}\angle C\right)=\sqrt{\frac{1-\frac{10}{17}}{2}}=\sqrt{\frac{7}{34}} \\ \sin\left(\frac{1}{2}\angle A\right)=\sqrt{\frac{1-\frac{43}{85}}{2}}=\sqrt{\frac{21}{85}}.\end{align*}
 Thus the answer is 
 \begin{align*} & 30\cdot 34\cdot\sqrt{\frac{3}{10}\cdot\frac{7}{34}\cdot\frac{21}{85}} \\ =&~30\cdot 34\cdot\sqrt{\frac{3}{2\cdot 5}\cdot\frac{7}{2\cdot 17}\cdot\frac{3\cdot 7}{5\cdot 17}} \\ =&~30\cdot 34\cdot\frac{3\cdot 7}{2\cdot 5\cdot 17} \\ =&~(2\cdot 3\cdot 5)\cdot(2\cdot 17)\cdot\frac{3\cdot 7}{2\cdot 5\cdot 17} \\ =&~2\cdot 3^{2}\cdot 7 \\ =&~\boxed{126}.\end{align*}





\textbf{Solution 3 (A lengthier, but less trigonometric approach)}

First, instead of using angles to find $[AI_1I_2]$, let's try to find the area of other, simpler figures, and subtract that from $[ABC]$. However, to do this, we need to be able to figure out the length of the inradii, and so, we need to find $AX$. 

 
To minimize $[AI_1I_2]$, intuitively, we should try to minimize the length of $AX$, since, after using the $rs=A$ formula for the area of a triangle, we'll be able to minimize the inradii lengths, and thus, eventually minimize the area of $[AI_1I_2]$. (Proof needed here).

 


 We need to minimize $AX$. Let $AX=d$, $BX=s$, and $CX=32-s$. After an application of Stewart's Theorem, we will get that 
 \[d=\sqrt{s^2-24s+900}\] To minimize this quadratic, $s=12$ whereby we conclude that $d=6\sqrt{21}$. 

 


 From here, draw perpendiculars down from $I_1$ and $I_2$ to $AB$ and $AC$ respectively, and label the foot of these perpendiculars $D$ and $E$ respectively. After, draw the inradii from $I_1$ to $BX$, and from $I_2$ to $CX$, and draw in $I_1I_2$. 

 Label the foot of the inradii to $BX$ and $CX$, $F$ and $G$, respectively. From here, we see that to find $[AI_1I_2]$, we need to find $[ABC]$, and subtract off the sum of $[DBCEI_2I_1], [ADI_1],$ and $[AEI_2]$. 

 
$[DBCEI_2I_1]$ can be found by finding the area of two quadrilaterals $[DBFI_1]+[ECGI_2]$ as well as the area of a trapezoid $[FGI_2I_1]$. If we let the inradius of $ABX$ be $r_1$ and if we let the inradius of $ACX$ be $r_2$, we'll find, after an application of basic geometry and careful calculations on paper, that $[DBCEI_2I_1]=13r_1+19r_2$. 

 
The area of two triangles can be found in a similar fashion, however, we must use $XYZ$ substitution to solve for $AD$ as well as $AE$. After doing this, we'll get a similar sum in terms of $r_1$ and $r_2$ for the area of those two triangles which is equal to 
 \[\frac{(9+3\sqrt{21})(r_1)}{2} + \frac{(7+3\sqrt{21})(r_2)}{2}.\]

 Now we're set. Summing up the area of the Hexagon and the two triangles and simplifying, we get that the formula for $[AI_1I_2]$ is just 
 \[[ABC]-\left(\frac{(35+3\sqrt{21})(r_1)}{2}+\frac{(45+3r_2\sqrt{21})(r_2)}{2}\right).\] 

 Using Heron's formula, $[ABC]=96\sqrt{21}$. Solving for $r_1$ and $r_2$ using Heron's in $ABX$ and $ACX$, we get that $r_1=3\sqrt{21}-9$ and $r_2=3\sqrt{21}-7$. From here, we just have to plug into our above equation and solve.

 Doing so gets us that the minimum area of $AI_1I_2=\boxed{126}.$

 -Azeem H.(Mathislife52) ~edited by phoenixfire

 


	\item (\textit{AIME I 2018}) Let $SP_1P_2P_3EP_4P_5$ be a heptagon. A frog starts jumping at vertex $S$. From any vertex of the heptagon except $E$, the frog may jump to either of the two adjacent vertices. When it reaches vertex $E$, the frog stops and stays there. Find the number of distinct sequences of jumps of no more than $12$ jumps that end at $E$.


 
 \textbf{Answer: }351



\textbf{Solution 1}

This is easily solved by recursion/dynamic programming. To simplify the problem somewhat, let us imagine the seven vertices on a line $E \leftrightarrow P_4 \leftrightarrow P_5 \leftrightarrow S \leftrightarrow P_1 \leftrightarrow P_2 \leftrightarrow P_3 \leftrightarrow E$. We can count the number of left/right (L/R) paths of length $\le 11$ that start at $S$ and end at either $P_4$ or $P_3$.

 We can visualize the paths using the common grid counting method by starting at the origin $(0,0)$, so that a right (R) move corresponds to moving 1 in the positive $x$ direction, and a left (L) move corresponds to moving 1 in the positive $y$ direction. Because we don't want to move more than 2 units left or more than 3 units right, our path must not cross the lines $y = x+2$ or $y = x-3$. Letting $p(x,y)$ be the number of such paths from $(0,0)$ to $(x,y)$ under these constraints, we have the following base cases:

 $p(x,0) = \begin{cases} 1 & x \le 3 \\ 0 & x > 3 \end{cases} \qquad p(0,y) = \begin{cases} 1 & y \le 2 \\ 0 & y > 2 \end{cases}$

 and recursive step$p(x,y) = p(x-1,y) + p(x,y-1)$ for $x,y \ge 1$.

 The filled in grid will look something like this, where the lower-left $1$ corresponds to the origin:

 $\begin{tabular}{|c|c|c|c|c|c|c|c|} \hline 0 & 0 & 0 & 0 & \textbf{89} & & & \\ \hline 0 & 0 & 0 & \textbf{28} & 89 & & & \\ \hline 0 & 0 & \textbf{9} & 28 & 61 & 108 & 155 & \textbf{155} \\ \hline 0 & \textbf{3} & 9 & 19 & 33 & 47 & \textbf{47} & 0 \\ \hline \textbf{1} & 3 & 6 & 10 & 14 & \textbf{14} & 0 & 0 \\ \hline 1 & 2 & 3 & 4 & \textbf{4} & 0 & 0 & 0 \\ \hline 1 & 1 & 1 & \textbf{1} & 0 & 0 & 0 & 0 \\ \hline \end{tabular}$

 The bolded numbers on the top diagonal represent the number of paths from $S$ to $P_4$ in 2, 4, 6, 8, 10 moves, and the numbers on the bottom diagonal represent the number of paths from $S$ to $P_3$ in 3, 5, 7, 9, 11 moves. We don't care about the blank entries or entries above the line $x+y = 11$. The total number of ways is $1+3+9+28+89+1+4+14+47+155 = \boxed{351}$.

 (Solution by scrabbler94)

 



\textbf{Solution 2}

Let $E_n$ denotes the number of sequences with length $n$ that ends at $E$. Define similarly for the other vertices. We seek for a recursive formula for $E_n$.
 \begin{align*} E_n&=P_{3_{n-1}}+P_{4_{n-1}} \\ &=P_{2_{n-2}}+P_{5_{n-2}} \\ &=P_{1_{n-3}}+P_{3_{n-3}}+S_{n-3}+P_{4_{n-3}} \\ &=(P_{3_{n-3}}+P_{4_{n-3}})+S_{n-3}+P_{1_{n-3}} \\ &=E_{n-2}+S_{n-3}+P_{1_{n-3}} \\ &=E_{n-2}+P_{1_{n-4}}+P_{5_{n-4}}+S_{n-4}+P_{2_{n-4}} \\ &=E_{n-2}+(S_{n-4}+P_{1_{n-4}})+P_{5_{n-4}}+P_{2_{n-4}} \\ &=E_{n-2}+(E_{n-1}-E_{n-3})+S_{n-5}+P_{4_{n-5}}+P_{1_{n-5}}+P_{3_{n-5}} \\ &=E_{n-2}+(E_{n-1}-E_{n-3})+(S_{n-5}+P_{1_{n-5}})+(P_{4_{n-5}}+P_{3_{n-5}}) \\ &=E_{n-2}+(E_{n-1}-E_{n-3})+(E_{n-2}-E_{n-4})+E_{n-4} \\ &=E_{n-1}+2E_{n-2}-E_{n-3} \\ \end{align*}
Computing a few terms we have $E_0=0$, $E_1=0$, $E_2=0$, $E_3=1$, and $E_4=1$.

 Using the formula yields $E_5=3$, $E_6=4$, $E_7=9$, $E_8=14$, $E_9=28$, $E_{10}=47$, $E_{11}=89$, and $E_{12}=155$.

 Finally adding yields $\sum_{k=0}^{12}E_k=\boxed{351}$.

 \textbf{Get solution faster in Sol 2} In the 5th to last line, we have $P_{5_{n-4}}+P_{2_{n-4}}=E_{n-2}$ by shifting the index in line 2.

 



\textbf{Solution 3}

For vertices $S, P_1, P_2, P_5,$ the number of ways to get there after $n$ jumps is the sum of the number of ways to get to the adjacent vertices after $n-1$ jumps. For vertices $P_3,$ and $P_4,$ the number of ways to get there after $n$ jumps is he number of ways to get to $P_2$ and $P_5$ after $n-1$ jumps.

 $\begin{tabular}{|l|c|c|c|c|c|c|c|} \hline Jump \# & \(S\)   & \(P_1\) & \(P_2\) & \(P_3\) & \(E\)        & \(P_4\) & \(P_5\) \\ \hline 1       & 0       & 1       & 0       & 0       & \textbf{0}   & 0       & 1       \\ \hline 2       & 2       & 0       & 1       & 0       & \textbf{0}   & 1       & 0       \\ \hline 3       & 0       & 3       & 0       & 1       & \textbf{1}   & 0       & 3       \\ \hline 4       & 6       & 0       & 4       & 0       & \textbf{1}   & 3       & 0       \\ \hline 5       & 0       & 10      & 0       & 4       & \textbf{3}   & 0       & 9       \\ \hline 6       & 19      & 0       & 14      & 0       & \textbf{4}   & 9       & 0       \\ \hline 7       & 0       & 33      & 0       & 14      & \textbf{9}   & 0       & 28      \\ \hline 8       & 61      & 0       & 47      & 0       & \textbf{14}  & 28      & 0       \\ \hline 9       & 0       & 108     & 0       & 47      & \textbf{28}  & 0       & 89      \\ \hline 10      & 197     & 0       & 155     & 0       & \textbf{47}  & 89      & 0       \\ \hline 11      & 0       & 352     & 0       & 155     & \textbf{89}  & 0       & 286     \\ \hline 12      & 638     & 0       & 507     & 0       & \textbf{155} & 286     & 0       \\ \hline \end{tabular}$
$0+0+1+1+3+4+9+14+28+47+89+155=\boxed{351}.$



	\item (\textit{AIME I 2018}) David found four sticks of different lengths that can be used to form three non-congruent convex cyclic quadrilaterals, $A,\text{ }B,\text{ }C$, which can each be inscribed in a circle with radius $1$. Let $\varphi_A$ denote the measure of the acute angle made by the diagonals of quadrilateral $A$, and define $\varphi_B$ and $\varphi_C$ similarly. Suppose that $\sin\varphi_A=\frac{2}{3}$, $\sin\varphi_B=\frac{3}{5}$, and $\sin\varphi_C=\frac{6}{7}$. All three quadrilaterals have the same area $K$, which can be written in the form $\dfrac{m}{n}$, where $m$ and $n$ are relatively prime positive integers. Find $m+n$.


 
 \textbf{Answer: }059



\textbf{Solution 1}

Suppose our four sides lengths cut out arc lengths of $2a$, $2b$, $2c$, and $2d$, where $a+b+c+d=180^\circ$. Then, we only have to consider which arc is opposite $2a$. These are our three cases, so
 \[\varphi_A=a+c\]
 \[\varphi_B=a+b\]
 \[\varphi_C=a+d\]Our first case involves quadrilateral $ABCD$ with $\overarc{AB}=2a$, $\overarc{BC}=2b$, $\overarc{CD}=2c$, and $\overarc{DA}=2d$.

 Then, by Law of Sines, $AC=2\sin\left(\frac{\overarc{ABC}}{2}\right)=2\sin(a+b)$ and $BD=2\sin\left(\frac{\overarc{BCD}}{2}\right)=2\sin(a+d)$. Therefore,

 
 \[K=\frac{1}{2}\cdot AC\cdot BD\cdot \sin(\varphi_A)=2\sin\varphi_A\sin\varphi_B\sin\varphi_C=\frac{24}{35},\]so our answer is $24+35=\boxed{059}$.

 Note that the conditions of the problem are satisfied when the lengths of the four sticks are about $0.32, 0.91, 1.06, 1.82$.

 By S.B.

 \textbf{Note} 
                    \begin{figure}[H]
                        \centering
                        \includegraphics[width=0.5\textwidth]{Solution/AIME_I_2018_15_Solution_Figure_1.png}
                    \end{figure}
                    

 The solution uses 
 \[\varphi_A=a+c.\]

 We can see that this follows because $\varphi_A = \frac12 (2a+2c)=a+c,$ where $a$ and $c$ are the central angles of opposite sides.\_\_\_\_Shen Kislay Kai

 



\textbf{Solution 2}

Suppose the four side lengths of the quadrilateral cut out arc lengths of $2a$, $2b$, $2c$, and $2d$. $a+b+c+d=180^\circ$. Therefore, without losing generality, 

 
 \[\varphi_A=a+b\]
 \[\varphi_B=b+c\]
 \[\varphi_C=a+c\]
$(1)+(3)-(2)$, $(1)+(2)-(3)$, and $(2)+(3)-(1)$ yields

 
 \[2a=\varphi_A+\varphi_C-\varphi_B\]
 \[2b=\varphi_A+\varphi_B-\varphi_C\]
 \[2c=\varphi_B+\varphi_C-\varphi_A\]

 Because $2d=360^\circ-2a-2b-2c,$Therefore,

 
 \[2d=360^\circ-\varphi_A-\varphi_B-\varphi_C\]

 Using the \href{https://artofproblemsolving.com/wiki/index.php/Trigonometric\_identities\#Sum-to-product\_identities}{sum-to-product identities}, our area of the quadrilateral $K$ then would be

 
 \begin{align*}  K&=\frac{1}{2}(\sin(2a)+\sin(2b)+\sin(2c)+\sin(2d))\\ &=\frac{1}{2}(\sin(\varphi_A+\varphi_B-\varphi_C)+\sin(\varphi_B+\varphi_C-\varphi_A)+\sin(\varphi_C+\varphi_A-\varphi_B)-\sin(\varphi_A+\varphi_B+\varphi_C))\\ &=\frac{1}{2}(2\sin\varphi_B\cos(\varphi_A-\varphi_C)-2\sin\varphi_B\cos(\varphi_A+\varphi_C))\\ &=\frac{1}{2}\cdot2\cdot2\sin\varphi_A\sin\varphi_B\sin\varphi_C\\ &=2\sin\varphi_A\sin\varphi_B\sin\varphi_C\\ &=\frac{24}{35}\\ \end{align*}


 Therefore, our answer is $24+35=\boxed{059}$.

 



\textbf{Solution 3}

Let the four stick lengths be $a$, $b$, $c$, and $d$. WLOG, let’s say that quadrilateral $A$ has sides $a$ and $d$ opposite each other, quadrilateral $B$ has sides $b$ and $d$ opposite each other, and quadrilateral $C$ has sides $c$ and $d$ opposite each other. The area of a convex quadrilateral can be written as $\frac{1}{2} d_1 d_2 \sin{\theta}$, where $d_1$ and $d_2$ are the lengths of the diagonals of the quadrilateral and $\theta$ is the angle formed by the intersection of $d_1$ and $d_2$. By Ptolemy's theorem $d_1 d_2 = ad+bc$ for quadrilateral $A$, so, defining $K_A$ as the area of $A$,
 \[K_A = \frac{1}{2} (ad+bc)\sin{\varphi_A}\]Similarly, for quadrilaterals $B$ and $C$,
 \[K_B = \frac{1}{2} (bd+ac)\sin{\varphi_B}\]and
 \[K_C = \frac{1}{2} (cd+ab)\sin{\varphi_C}\]Multiplying the three equations and rearranging, we see that
 \[K_A K_B K_C = \frac{1}{8} (ab+cd)(ac+bd)(ad+bc)\sin{\varphi_A}\sin{\varphi_B}\sin{\varphi_B}\]
 \[K^3 = \frac{1}{8} (ab+cd)(ac+bd)(ad+bc)\left(\frac{2}{3}\right) \left(\frac{3}{5}\right) \left(\frac{6}{7}\right)\]
 \[\frac{70}{3}K^3 = (ab+cd)(ac+bd)(ad+bc)\]The circumradius $R$ of a cyclic quadrilateral with side lengths $a$, $b$, $c$, and $d$ and area $K$ can be computed as $R = \frac{\sqrt{(ab+cd)(ac+bd)(ad+bc)}}{4K}$.Inserting what we know,
 \[1 = \frac{\sqrt{\frac{70}{3}K^3}}{4K}\quad \Rightarrow \quad 16K^2 = \frac{70}{3}K^3\quad \Rightarrow \quad \frac{24}{35} = K\]So our answer is $24 + 35 = \boxed{059}$.

 



\textbf{Solution 4 (No words)}


                    \begin{figure}[H]
                        \centering
                        \includegraphics[width=0.5\textwidth]{Solution/AIME_I_2018_15_Solution_Figure_2.png}
                    \end{figure}
                    

 \textbf{vladimir.shelomovskii@gmail.com, vvsss}

 



\textbf{Solution 5}

Let the sides of the quadrilaterals be $a,b,c,$ and $d$ in some order such that $A$ has $a$ opposite of $c$, $B$ has $a$ opposite of $b$, and $C$ has $a$ opposite of $d$. Then, let the diagonals of $A$ be $e$ and $f$. Similarly to solution $2$, we get that $\frac{2}{3}(ac+bd)=\frac{3}{5}(ab+cd)=\frac{6}{7}(ad+bc)=2K$, but this is also equal to $2\cdot\frac{eab+ecd}{4(1)}=2\cdot\frac{fad+fbc}{4(1)}$ using the area formula for a triangle using the circumradius and the sides, so $\frac{e(ab+cd)}{2}=\frac{3}{5}(ab+cd)$ and $\frac{f(ad+bc)}{2}=\frac{6}{7}(ad+bc)$. Solving for $e$ and $f$, we get that $e=\frac{6}{5}$ and $f=\frac{12}{7}$, but $K=\frac{1}{2}\cdot\frac{2}{3}\cdot{}ef$, similarly to solution $2$, so $K=\frac{24}{35}$ and the answer is $24+35=\boxed{059}$.

 


\end{enumerate}